% How accurate LIGHT_Kr v0_5 is: the pure-krypton model, scored against
% pure-krypton measurements, with its propagated uncertainty, its
% reference point and its checks. No number here is typed by hand except
% the measurements, which carry their sources. They come from
%   overview_v0_5_numbers.tex     make_overview_v0_5.py     (\wg)
%   system_map_v0_5_numbers.tex   make_system_map_v0_5.py   (\wm)
% Build from LIGHT_Kr/docs:
%   python3 make_system_map_v0_5.py && python3 make_overview_v0_5.py
%   latexmk -pdf accuracy_overview_v0_5.tex
\documentclass[11pt]{article}
\usepackage[a4paper,margin=2.2cm]{geometry}
\usepackage[T1]{fontenc}
\usepackage[scaled=0.95]{helvet}
\renewcommand{\familydefault}{\sfdefault}
\usepackage{amsmath,amssymb}
\usepackage{sansmath}
\usepackage{booktabs,tabularx,array}
\usepackage{xcolor}
\usepackage{float}
\usepackage{enumitem}
\usepackage{graphicx}
\usepackage{tikz}
\usepackage{pgfplots}
\pgfplotsset{compat=1.18}
\usetikzlibrary{arrows.meta,decorations.pathmorphing}
\setlist{nosep, leftmargin=1.2em}
\setlength{\parindent}{0pt}
\setlength{\parskip}{5pt}
\definecolor{ink}{HTML}{0B0B0B}
\definecolor{ink2}{HTML}{52514E}
\definecolor{cE}{HTML}{2A78D6}
\definecolor{cX}{HTML}{EDA100}
\definecolor{cP}{HTML}{C46A00}
\definecolor{cV}{HTML}{A63A83}
\definecolor{cH}{HTML}{E34948}
\definecolor{cI}{HTML}{52514E}
\renewcommand{\topfraction}{0.95}
\renewcommand{\textfraction}{0.03}

% LIGHT_Kr v0_5 map numbers. GENERATED by make_system_map_v0_5.py
% from the model (0.5.0-kr). Do not edit. Rerun the script instead.
\providecommand\wmdef[2]{\expandafter\def\csname wm@#1\endcsname{#2}}
\wmdef{e.P}{100.0}
\wmdef{e.pen}{1.0}
\wmdef{e.pen2}{0.0}
\wmdef{e.super}{0.8}
\wmdef{e.eq}{0.0}
\wmdef{e.G1}{8.9}
\wmdef{e.C1}{0.5}
\wmdef{e.HM}{0.0}
\wmdef{e.C2}{0.0}
\wmdef{e.76}{73.3}
\wmdef{e.stepS}{4.0}
\wmdef{e.stepP}{1.4}
\wmdef{e.ei}{5.3}
\wmdef{e.emix}{0.0}
\wmdef{e.em}{0.0}
\wmdef{e.G6}{3.2}
\wmdef{e.dU}{0.0}
\wmdef{e.cap}{5.2}
\wmdef{G1.krs}{8.9}
\wmdef{G1.heat}{0.0}
\wmdef{C1.kr}{0.5}
\wmdef{C1.heat}{0.0}
\wmdef{HM.ion}{0.0}
\wmdef{HM.heat}{0.0}
\wmdef{kpp.back}{28.1}
\wmdef{K76.heat}{18.8}
\wmdef{KPS.heat}{57.3}
\wmdef{PRAD.hv}{0.5}
\wmdef{PRAD.heat}{0.0}
\wmdef{kpp.stepP}{5.8}
\wmdef{stepP.heat}{0.0}
\wmdef{kpp.left}{0.0}
\wmdef{krs.G2}{3.2}
\wmdef{krs.AK8}{21.8}
\wmdef{krs.AK8f}{25.4}
\wmdef{krs.AK8r}{3.5}
\wmdef{krs.stepS}{11.0}
\wmdef{krs.super}{0.8}
\wmdef{krs.emix}{0.0}
\wmdef{krs.res}{0.0}
\wmdef{krs.mix}{0.0}
\wmdef{MIX.heat}{0.0}
\wmdef{SUPER.heat}{0.0}
\wmdef{EMIX.heat}{0.0}
\wmdef{RES.hv}{0.0}
\wmdef{krs.left}{0.1}
\wmdef{G2.ion}{2.2}
\wmdef{C2.ion}{0.0}
\wmdef{stepS.ion}{15.1}
\wmdef{stepS.heat}{0.0}
\wmdef{stepP.ion}{7.2}
\wmdef{AK8.heat}{1.8}
\wmdef{k2s.RAD}{6.6}
\wmdef{k2s.Q}{0.0}
\wmdef{k2s.pen}{0.0}
\wmdef{k2s.ei}{13.3}
\wmdef{k2s.eq}{0.0}
\wmdef{k2s.em}{0.0}
\wmdef{k2s.left}{0.1}
\wmdef{RAD.hv}{6.1}
\wmdef{RAD.heat}{0.6}
\wmdef{Q.heat}{0.0}
\wmdef{pen2.ion}{0.0}
\wmdef{ei.ion}{18.5}
\wmdef{ei.heat}{0.0}
\wmdef{eq.heat}{0.0}
\wmdef{em.heat}{0.0}
\wmdef{K15.heat}{2.1}
\wmdef{R11p.in}{41.6}
\wmdef{R11p.kr}{36.6}
\wmdef{R11p.heat}{5.0}
\wmdef{ions.left}{-0.7}
\wmdef{out.147}{6.1}
\wmdef{out.IR}{0.5}
\wmdef{out.res}{0.0}
\wmdef{out.heat}{90.7}
\wmdef{out.elastic}{3.2}
\wmdef{out.ions}{-0.7}
\wmdef{out.other}{0.3}
\wmdef{e.ion}{10.7}
\wmdef{rung.heat}{76.1}
\wmdef{eion.ion}{10.7}
\wmdef{fac.x1s}{0.89}
\wmdef{fac.x2p}{0.92}
\wmdef{fac.xhi}{0.75}
\wmdef{fac.mlo}{0.93}
\wmdef{fac.mhi}{1.23}
\wmdef{version}{0.5.0-kr}
\wmdef{md5}{b392d33a7c61}
\wmdef{p}{1}
\wmdef{Tg0}{300}
\wmdef{P}{300}
\wmdef{pulse}{1}
\wmdef{tend}{2}
\wmdef{ne}{10^{15}}
\wmdef{N}{2.45\times10^{19}}
\wmdef{Edep}{0.30}
\wmdef{Te}{1.57}
\wmdef{TeMax}{2.11}
\wmdef{enLo}{1.3}
\wmdef{enHi}{6.44}
\wmdef{kinN}{338}
\wmdef{kinOn}{263}
\wmdef{kinRatioLo}{1.0001}
\wmdef{kinRatioHi}{1.0002}
\wmdef{kinGap}{10^{-4}}
\wmdef{closure}{4\times10^{-15}}
\wmdef{TgEnd}{856}
\wmdef{TgMax}{856}
\wmdef{pEnd}{2.85}
\wmdef{nePeak}{1.00\times10^{15}}
\wmdef{neEnd}{1.14\times10^{13}}
\wmdef{krsFormed}{1.09\times10^{18}}
\wmdef{kppFormed}{1.07\times10^{18}}
\wmdef{k2sFormed}{4.75\times10^{16}}
\wmdef{brs.76}{93.1}
\wmdef{brs.AK8}{4.4}
\wmdef{brs.G2}{0.5}
\wmdef{brs.step}{1.9}
\wmdef{brs.super}{0.1}
\wmdef{brp.KPS}{89.0}
\wmdef{brp.76r}{9.5}
\wmdef{brp.PRAD}{0.6}
\wmdef{brp.step}{0.9}
\wmdef{brk2.RAD}{28.5}
\wmdef{brk2.redis}{14.0}
\wmdef{brk2.ei}{56.9}
\wmdef{k2s.leftpc}{0.6}
\wmdef{gs.G1}{1.5}
\wmdef{gs.KPS}{87.9}
\wmdef{gs.PRAD}{0.6}
\wmdef{gs.76r}{9.4}
\wmdef{gs.redis}{0.6}
\wmdef{form.1s5}{43.9}
\wmdef{form.1s4}{50.6}
\wmdef{form.1s3}{5.5}
\wmdef{redis.pc}{14.0}
\wmdef{cyc.perK2}{25}
\wmdef{cyc.e}{73.3}
\wmdef{cyc.net}{9.10\times10^{17}}
\wmdef{ion.G2}{5}
\wmdef{ion.C2}{0.089}
\wmdef{ion.stepS}{34}
\wmdef{ion.stepP}{16}
\wmdef{ion.ei}{45}
\wmdef{k.G2}{5.0\times10^{-10}}
\wmdef{k.AK8}{3.8\times10^{-32}}
\wmdef{k.AK8.1s4}{9.2\times10^{-32}}
\wmdef{tpow.AK8}{-0.64}
\wmdef{tauRAD}{270}
\wmdef{tauRADF}{3.4}
\wmdef{hv.RAD}{8.43}
\wmdef{E.kr2s}{9.21}
\wmdef{E.kr2p}{12.851}
\wmdef{k.K15}{2.4\times10^{-31}}
\wmdef{tpow.K15}{-0.75}
\wmdef{k.R11p}{1.6\times10^{-6}}
\wmdef{pow.R11p}{0.55}
\wmdef{k.KPS}{1.7\times10^{-10}}
\wmdef{k.MIX}{1.3\times10^{-13}}
\wmdef{tauPRAD}{37.9}
\wmdef{nLines}{28}
\wmdef{nPairs}{26}
\wmdef{E.ei}{3.64}
\wmdef{E.stepS}{4.08}
\wmdef{E.stepP}{2.70}
\wmdef{E.PP}{11.30}
\wmdef{E.S}{9.9152}
\wmdef{E.1s5}{9.9152}
\wmdef{E.1s4}{10.0324}
\wmdef{E.1s3}{10.5624}
\wmdef{E.1s2}{10.6436}
\wmdef{E.ion}{13.9996}
\wmdef{E.ai}{12.85}
\wmdef{nG1}{4}
\wmdef{nC1}{47}
\wmdef{nHM}{20}
\wmdef{C1.lo}{11.30}
\wmdef{C1.hi}{13.6}
\wmdef{G1.hi}{10.64}
\wmdef{gap5p}{1.388}
\wmdef{pumpx}{18.8}
\wmdef{gap5p.lo}{0.660}
\wmdef{pen.e}{5.831}
\wmdef{sig.76}{1.1\times10^{-14}}
\wmdef{nodes}{30}
\wmdef{nrows}{219}
\wmdef{notes.krs}{9.4}
\wmdef{notes.krp}{43.0}
\wmdef{notes.k2s}{21.4}
\wmdef{om.kannari}{6.71}
\wmdef{om.nocycle}{11.42}
\wmdef{om.twot}{5.58}
\wmdef{om.twot.te}{1.53}
\wmdef{bd.ei_kr2.def}{1}
\wmdef{bd.ei_kr2.lo}{0}
\wmdef{bd.ei_kr2.hi}{2}
\wmdef{bd.ei_kr2.lo147}{15.43}
\wmdef{bd.ei_kr2.hi147}{4.47}
\wmdef{bd.ei_kr2.cell}{15.43 to 4.47}
\wmdef{bd.kstep_s.def}{1}
\wmdef{bd.kstep_s.lo}{0}
\wmdef{bd.kstep_s.hi}{2}
\wmdef{bd.kstep_s.lo147}{6.96}
\wmdef{bd.kstep_s.hi147}{5.60}
\wmdef{bd.kstep_s.cell}{6.96 to 5.60}
\wmdef{bd.kstep_p.def}{1}
\wmdef{bd.kstep_p.lo}{0}
\wmdef{bd.kstep_p.hi}{2}
\wmdef{bd.kstep_p.lo147}{6.37}
\wmdef{bd.kstep_p.hi147}{5.83}
\wmdef{bd.kstep_p.cell}{6.37 to 5.83}
\wmdef{bd.k_5s5p.def}{1}
\wmdef{bd.k_5s5p.lo}{0.6}
\wmdef{bd.k_5s5p.hi}{1.4}
\wmdef{bd.k_5s5p.lo147}{7.06}
\wmdef{bd.k_5s5p.hi147}{5.46}
\wmdef{bd.k_5s5p.cell}{7.06 to 5.46}
\wmdef{bd.k_5s5p_j1.def}{1}
\wmdef{bd.k_5s5p_j1.lo}{0}
\wmdef{bd.k_5s5p_j1.hi}{1}
\wmdef{bd.k_5s5p_j1.lo147}{5.55}
\wmdef{bd.k_5s5p_j1.hi147}{6.08}
\wmdef{bd.k_5s5p_j1.cell}{5.55 to \textbf{6.08}}
\wmdef{bd.kps_kr.def}{1.7\times10^{-10}}
\wmdef{bd.kps_kr.lo}{3.0\times10^{-11}}
\wmdef{bd.kps_kr.hi}{1.7\times10^{-10}}
\wmdef{bd.kps_kr.lo147}{7.11}
\wmdef{bd.kps_kr.hi147}{6.08}
\wmdef{bd.kps_kr.cell}{7.11 to \textbf{6.08}}
\wmdef{bd.kps_5s.def}{1}
\wmdef{bd.kps_5s.lo}{0}
\wmdef{bd.kps_5s.hi}{1}
\wmdef{bd.kps_5s.lo147}{6.36}
\wmdef{bd.kps_5s.hi147}{6.08}
\wmdef{bd.kps_5s.cell}{6.36 to \textbf{6.08}}
\wmdef{bd.kr_pen.def}{5.0\times10^{-10}}
\wmdef{bd.kr_pen.lo}{1.5\times10^{-10}}
\wmdef{bd.kr_pen.hi}{5.0\times10^{-10}}
\wmdef{bd.kr_pen.lo147}{6.50}
\wmdef{bd.kr_pen.hi147}{6.08}
\wmdef{bd.kr_pen.cell}{6.50 to \textbf{6.08}}
\wmdef{bd.kr2_form.def}{3.8\times10^{-32}}
\wmdef{bd.kr2_form.lo}{1.7\times10^{-32}}
\wmdef{bd.kr2_form.hi}{4.2\times10^{-32}}
\wmdef{bd.kr2_form.lo147}{5.34}
\wmdef{bd.kr2_form.hi147}{6.20}
\wmdef{bd.kr2_form.cell}{5.34 to 6.20}
\wmdef{bd.kr2_form_1s4.def}{9.2\times10^{-32}}
\wmdef{bd.kr2_form_1s4.lo}{1.3\times10^{-32}}
\wmdef{bd.kr2_form_1s4.hi}{9.2\times10^{-32}}
\wmdef{bd.kr2_form_1s4.lo147}{4.29}
\wmdef{bd.kr2_form_1s4.hi147}{6.08}
\wmdef{bd.kr2_form_1s4.cell}{4.29 to \textbf{6.08}}
\wmdef{bd.kr2_form_1s3.def}{3.8\times10^{-32}}
\wmdef{bd.kr2_form_1s3.lo}{0}
\wmdef{bd.kr2_form_1s3.hi}{9.2\times10^{-32}}
\wmdef{bd.kr2_form_1s3.lo147}{5.81}
\wmdef{bd.kr2_form_1s3.hi147}{6.34}
\wmdef{bd.kr2_form_1s3.cell}{5.81 to 6.34}
\wmdef{bd.kr2_fast.def}{0}
\wmdef{bd.kr2_fast.lo}{0}
\wmdef{bd.kr2_fast.hi}{0.17}
\wmdef{bd.kr2_fast.lo147}{6.08}
\wmdef{bd.kr2_fast.hi147}{6.95}
\wmdef{bd.kr2_fast.cell}{\textbf{6.08} to 6.95}
\wmdef{bd.kr2_quench.def}{0}
\wmdef{bd.kr2_quench.lo}{0}
\wmdef{bd.kr2_quench.hi}{8.6\times10^{-15}}
\wmdef{bd.kr2_quench.lo147}{6.08}
\wmdef{bd.kr2_quench.hi147}{5.88}
\wmdef{bd.kr2_quench.cell}{\textbf{6.08} to 5.88}
\wmdef{bd.kr2_d0.def}{0.709}
\wmdef{bd.kr2_d0.lo}{0.62}
\wmdef{bd.kr2_d0.hi}{0.709}
\wmdef{bd.kr2_d0.lo147}{5.35}
\wmdef{bd.kr2_d0.hi147}{6.08}
\wmdef{bd.kr2_d0.cell}{5.35 to \textbf{6.08}}
\wmdef{bd.kr2_pen.def}{0}
\wmdef{bd.kr2_pen.lo}{0}
\wmdef{bd.kr2_pen.hi}{5.0\times10^{-10}}
\wmdef{bd.kr2_pen.lo147}{6.08}
\wmdef{bd.kr2_pen.hi147}{5.12}
\wmdef{bd.kr2_pen.cell}{\textbf{6.08} to 5.12}
\wmdef{bd.eq_kr2.def}{0}
\wmdef{bd.eq_kr2.lo}{0}
\wmdef{bd.eq_kr2.hi}{1.0\times10^{-7}}
\wmdef{bd.eq_kr2.lo147}{6.08}
\wmdef{bd.eq_kr2.hi147}{1.27}
\wmdef{bd.eq_kr2.cell}{\textbf{6.08} to 1.27}
\wmdef{bd.em_kr2.def}{0}
\wmdef{bd.em_kr2.lo}{0}
\wmdef{bd.em_kr2.hi}{1.0\times10^{-7}}
\wmdef{bd.em_kr2.lo147}{6.08}
\wmdef{bd.em_kr2.hi147}{12.44}
\wmdef{bd.em_kr2.cell}{\textbf{6.08} to 12.44}
\wmdef{bd.emix_5s.def}{4.3\times10^{-8}}
\wmdef{bd.emix_5s.lo}{4.3\times10^{-8}}
\wmdef{bd.emix_5s.hi}{2.3\times10^{-7}}
\wmdef{bd.emix_5s.lo147}{6.08}
\wmdef{bd.emix_5s.hi147}{6.10}
\wmdef{bd.emix_5s.cell}{\textbf{6.08} to 6.10}
\wmdef{bd.dr_tg.def}{0}
\wmdef{bd.dr_tg.lo}{0}
\wmdef{bd.dr_tg.hi}{1}
\wmdef{bd.dr_tg.lo147}{6.08}
\wmdef{bd.dr_tg.hi147}{6.50}
\wmdef{bd.dr_tg.cell}{\textbf{6.08} to 6.50}
\wmdef{bd.res_escape.def}{0}
\wmdef{bd.res_escape.lo}{0}
\wmdef{bd.res_escape.hi}{1.7\times10^{5}}
\wmdef{bd.res_escape.lo147}{6.08}
\wmdef{bd.res_escape.hi147}{6.05}
\wmdef{bd.res_escape.cell}{\textbf{6.08} to 6.05}
\wmdef{bd.hm_kr.def}{0}
\wmdef{bd.hm_kr.lo}{0}
\wmdef{bd.hm_kr.hi}{1}
\wmdef{bd.hm_kr.lo147}{6.08}
\wmdef{bd.hm_kr.hi147}{6.08}
\wmdef{bd.hm_kr.cell}{\textbf{6.08} to 6.08}
\wmdef{bd.n}{21}
\wmdef{bd.nzero}{4}
\wmdef{bd.noff}{4}
\wmdef{bd.ncalc}{3}
\def\wmBands{\el{} + \Krts{} \ra \Krtp{} + 2\el{}, a scale on Vriens \& Smeets & $1$ & $0$ to $2$ & 15.43 to 4.47 & 15.54 to 4.06 & 0.43 to 0.49 & 801 to 866\\
\el{} + \Krs{} \ra \Krp{} + 2\el{}, a scale on Vriens \& Smeets & $1$ & $0$ to $2$ & 6.96 to 5.60 & 6.60 to 5.15 & 0.49 to 0.47 & 851 to 859\\
\el{} + \Krpp{} \ra \Krp{} + 2\el{}, a scale on Vriens \& Smeets & $1$ & $0$ to $2$ & 6.37 to 5.83 & 5.86 to 5.36 & 0.49 to 0.47 & 854 to 858\\
\el{} + \Krs{} \ra \el{} + \Krpp{}, a scale on the measured cross sections & $1$ & $0.6$ to $1.4$ & 7.06 to 5.46 & 6.32 to 5.12 & 0.45 to 0.50 & 850 to 860\\
the same, a scale on the $J=1$ levels' proxies & $1$ & $0$ to $1$ & 5.55 to 6.08 & 4.76 to 5.58 & 0.32 to 0.48 & 858 to 856\\
\Krpp{} + Kr \ra \Krs{} + Kr & $1.7\times10^{-10}$ & $3.0\times10^{-11}$ to $1.7\times10^{-10}$\unit{cm^3\,s^{-1}} & 7.11 to 6.08 & 6.39 to 5.58 & 2.11 to 0.48 & 840 to 856\\
the share of \Krpp{} + Kr's products over the 5s levels & $1$ & $0$ to $1$ & 6.36 to 6.08 & 5.89 to 5.58 & 0.47 to 0.48 & 855 to 856\\
\Krs{} + \Krs{} \ra \Krp{} + Kr + \el{} & $5.0\times10^{-10}$ & $1.5\times10^{-10}$ to $5.0\times10^{-10}$\unit{cm^3\,s^{-1}} & 6.50 to 6.08 & 5.98 to 5.58 & 0.48 to 0.48 & 853 to 856\\
\Krs(1s$_5$) + 2Kr \ra \Krts{} + Kr & $3.8\times10^{-32}$ & $1.7\times10^{-32}$ to $4.2\times10^{-32}$\unit{cm^6\,s^{-1}} & 5.34 to 6.20 & 4.88 to 5.70 & 0.49 to 0.48 & 860 to 855\\
\Krs(1s$_4$) + 2Kr \ra \Krts{} + Kr & $9.2\times10^{-32}$ & $1.3\times10^{-32}$ to $9.2\times10^{-32}$\unit{cm^6\,s^{-1}} & 4.29 to 6.08 & 3.89 to 5.58 & 0.51 to 0.48 & 866 to 856\\
\Krs(1s$_3$) + 2Kr \ra \Krts{} + Kr & $3.8\times10^{-32}$ & $0$ to $9.2\times10^{-32}$\unit{cm^6\,s^{-1}} & 5.81 to 6.34 & 5.33 to 5.83 & 0.48 to 0.48 & 857 to 855\\
the share of the 1s$_4$ formation into the fast 0$_u^+$ state & $0$ & $0$ to $0.17$ & 6.08 to 6.95 & 5.58 to 6.51 & 0.48 to 0.48 & 856 to 851\\
\Krts{} + Kr \ra 3Kr & $0$ & $0$ to $8.6\times10^{-15}$\unit{cm^3\,s^{-1}} & 6.08 to 5.88 & 5.58 to 5.41 & 0.48 to 0.48 & 856 to 857\\
the excimer's well, for its redissociation & $0.709$ & $0.62$ to $0.709$\unit{eV} & 5.35 to 6.08 & 4.90 to 5.58 & 0.49 to 0.48 & 859 to 856\\
\Krts{} + \Krts{} \ra \Krtp{} + 2Kr + \el{} & $0$ & $0$ to $5.0\times10^{-10}$\unit{cm^3\,s^{-1}} & 6.08 to 5.12 & 5.58 to 4.79 & 0.48 to 0.49 & 856 to 862\\
\el{} + \Krts{} \ra \el{} + 2Kr & $0$ & $0$ to $1.0\times10^{-7}$\unit{cm^3\,s^{-1}} & 6.08 to 1.27 & 5.58 to 0.94 & 0.48 to 0.52 & 856 to 885\\
\el{} + \Krts{} \ra \el{} + \Krts(0$_u^+$), pre-exponential & $0$ & $0$ to $1.0\times10^{-7}$\unit{cm^3\,s^{-1}} & 6.08 to 12.44 & 5.58 to 12.08 & 0.48 to 0.44 & 856 to 820\\
\el{} + \Krs(1s$_5$) \ra \el{} + \Krs(1s$_4$), pre-exponential & $4.3\times10^{-8}$ & $4.3\times10^{-8}$ to $2.3\times10^{-7}$\unit{cm^3\,s^{-1}} & 6.08 to 6.10 & 5.58 to 5.60 & 0.48 to 0.48 & 856 to 856\\
the gas-temperature factor on \Krtp{} + \el{} & $0$ & $0$ to $1$ & 6.08 to 6.50 & 5.58 to 5.99 & 0.48 to 0.47 & 856 to 854\\
the escaping resonance lines, per 1s$_4$ and 1s$_2$ & $0$ & $0$ to $1.7\times10^{5}$\unit{s^{-1}} & 6.08 to 6.05 & 5.58 to 5.55 & 0.48 to 0.48 & 856 to 856\\
Kr(nl) + Kr \ra \Krtp{} + \el{}, the share of the levels' excitations & $0$ & $0$ to $1$ & 6.08 to 6.08 & 5.58 to 5.58 & 0.48 to 0.48 & 856 to 856\\}

\input{overview_v0_5_numbers}
\makeatletter
\newcommand\wm[1]{\ifcsname wm@#1\endcsname\csname wm@#1\endcsname
  \else\errmessage{no map number '#1'. Rerun make_system_map_v0_5.py}\fi}
\newcommand\wg[1]{\ifcsname wg@#1\endcsname\csname wg@#1\endcsname
  \else\errmessage{no overview number '#1'. Rerun make_overview_v0_5.py}\fi}
\makeatother

\newcommand{\unit}[1]{\ensuremath{\,\mathrm{#1}}}
\newcommand{\el}{\ensuremath{\mathrm{e}^-}}
\newcommand{\Krs}{Kr\ensuremath{^*}}
\newcommand{\Krpp}{Kr\ensuremath{^{**}}}
\newcommand{\Krp}{Kr\ensuremath{^+}}
\newcommand{\Krtp}{Kr\ensuremath{_2^+}}
\newcommand{\Krts}{Kr\ensuremath{_2^*}}
\newcommand{\lv}[1]{1s\ensuremath{_#1}}
\newcommand{\ra}{\ensuremath{\to}\ }
\newcommand{\Edep}{\ensuremath{E_\mathrm{dep}}}
\newcommand{\x}[1]{\ensuremath{\times}#1}
\newcommand{\cm}[1]{\ensuremath{\,\mathrm{cm^{#1}s^{-1}}}}
% a pgfplots curve from a generated macro of coordinates
\newcommand\plotcoords[2]{\edef\tmp{\noexpand\addplot[#1] coordinates
  {#2};}\tmp}


\begin{document}
\sansmath

{\LARGE\bfseries How accurate LIGHT\_Kr v0\_5 is}\par
\vspace{4pt}
{\large The pure-krypton model LIGHT\_Kr v0\_5, version \wg{version}, on
16 October 2026: v0\_4 with the electrons' collisions with excited
krypton and with the excimer taken in, calculated where nothing is
measured, each 5s \ra 5p pumping charged its own energy, and the
two-temperature closure's tail table given its excited atoms, scored
against pure-krypton measurements, with its propagated uncertainty, its
reference point and its checks}\par
\vspace{10pt}

\noindent\fbox{\parbox{\dimexpr\linewidth-2\fboxsep-2\fboxrule\relax}{%
\setlength{\parskip}{5pt}%
\textbf{Bottom line.} No measurement in pure krypton covers v0\_5's own
regime, a gas at 0.1 to 2\unit{atm} whose electrons a field heats. The
measurements that exist test its parts. Only electron swarm measurements
train anything, the cross sections, which v0\_5 takes unchanged from
v0\_3 and v0\_4; every other measurement only scores the model.
\begin{itemize}
\item \textbf{What v0\_5 changes.} v0\_4 held at 0 the electrons'
collisions that ionize the excited atoms, and carried none with the
excimer \Krts. Where nothing is measured v0\_5 takes a published
calculation: the ionization of the 5s levels, of the 5p rung and of
\Krts{} by Vriens and Smeets's formula~\cite{vriens80}, which for \Krts{}
gives the $10^{-15}\unit{cm^2}$ of Flannery et al.'s calculation~\cite{flannery81}.
It charges each 5s \ra 5p pumping the energy of the 5p level it reaches
(v0\_4 charged the gap to the rung's lowest level), returns the 5s levels
to the ground state by detailed balance, mixes \lv{5} and \lv{4} by
electrons at Blagoev and Pilosov's measured rate~\cite{blagoev81}, and
gives the two-temperature closure's tail table an axis for the excited
atoms.
\item \textbf{The answer falls by a factor of \wg{chg.kin.ratio}.} At
LIGHT's reference point (\wm{p}\unit{atm}, \wm{P}\unit{kW\,cm^{-3}} for
\wm{pulse}\unit{\mu s}) v0\_5 sends \wg{ref.kin.147}\,\% of the input to
147\unit{nm} on the kinetic EEDF; v0\_4 sent \wg{chg.kin.v04}\,\%. The
electrons ionize \wg{ref.kin.k2ei}\,\% of the excimers before they
radiate (\wg{ref.kin.k2rad}\,\% do), and the stepwise ionizations and the
excimer's make most of the ions, which recombine onto the 5p rung, whose
relaxation heats the gas. With the three calculated rows at v0\_4's 0 the
share is \wg{undo.no_calc.kin}\,\%.
\item \textbf{Propagated from its sources}, the share carries
$\pm$\wg{unc.kin.quad} points in quadrature ($\pm$\wg{unc.kin.quad.pct}\,\%),
and the trained cross sections' posterior $\pm$\wg{unc.kin.xs.sd}
(Section~\ref{sec:accuracy}). The three largest terms are the electrons'
collisions with \Krts: its ionization, a calculation (the share
\wg{unc.end.ei_kr2.low}\,\% without it, \wg{unc.end.ei_kr2.high}\,\% at
twice it), and its quenching and the mixing of its two states, which
nothing measures or calculates (\wg{unc.end.eq_kr2.high} and
\wg{unc.end.em_kr2.high}\,\% at Kannari et al.'s estimate).
\item \textbf{A new score.} Blagoev et al.\ measured the 5s levels'
return to the ground state by thermal electrons in an
afterglow~\cite{blagoev84}. v0\_5's, the detailed balance of v0\_3's
trained excitation cross sections, lies at \wg{blag.1s5.lo} to
\wg{blag.1s5.hi} times theirs for \lv{5}, \wg{blag.1s4.lo} to
\wg{blag.1s4.hi} for \lv{4} and \wg{blag.1s3.lo} to \wg{blag.1s3.hi}
for \lv{3} (Section~\ref{sec:excited}); the return moves the reference
point by \wg{undo.no_super.kin.back}\,\%.
\item \textbf{The electroluminescence yield} of Mano et al.~\cite{mano22}
and \textbf{the electron swarm} are v0\_4's: \wg{el.lo} to \wg{el.hi} of
the yield's fit from 6 to 13.5\unit{Td}; the drift velocity within
\wg{sw.wtr.tr.rms}\,\% rms and the ionization coefficient within
\wg{sw.atr.tr.rms}\,\% on the train sets (Sections~\ref{sec:el}
and~\ref{sec:swarm}).
\item \textbf{Under an electron beam} the new rows act only where the
beam's plasma is dense: v0\_5 turns \wg{moz.300.kin} of the deposited
energy into 147-nm light at low pump, v0\_4's value, where Morozov et
al.\ measured \wg{moz.meas} $\pm$ \wg{moz.err}~\cite{morozov08}
(\wg{moz.kin.dev}\,\%), and at $3\times10^5\unit{W\,cm^{-3}}$ and
1\unit{atm} it is \wg{eck.kin.dev}\,\% from Eckstrom et al.'s report
(v0\_4: \wg{eck.kin.dev.v04}\,\%; Section~\ref{sec:beam}).
\item \textbf{The two-temperature closure} on its tail table with the
excited axis gives \wg{ref.2t.147}\,\% at the reference point,
\wg{tail.ref.f.dev}\,\% from the kinetic EEDF; on the table without the
axis (v0\_4's) \wg{tail.ref.f0}\,\%, \wg{tail.ref.f0.dev}\,\%. Its tail
now follows the kinetic one, and its Maxwellian bulk is what is left
(Section~\ref{sec:twot}).
\item \textbf{The numerical error is far smaller.} The energy ledger
closes to $\wg{ref.kin.closure}$ of the input; halving the kinetic
closure's renewal steps or doubling its cells moves the 147-nm and
near-IR shares by at most $\wg{conv.worst}$ of themselves.
\end{itemize}}}
\medskip

Table~\ref{tab:score} lists every test, Section~\ref{sec:accuracy} reads
them as errors and propagates the model's own uncertainty, and the
sections after it give each test's detail. The swarm datasets in the role
\emph{train} set the cross sections' factors (Section~\ref{sec:swarm});
every other measurement in the table only scores the model, and no
measurement set any other rate, which come from the sources in
Appendix~\ref{app:rates}.

\begin{table}[H]
\centering\footnotesize
\renewcommand{\arraystretch}{1.12}
\begin{tabularx}{\linewidth}{@{}>{\raggedright\arraybackslash}X
  >{\raggedright\arraybackslash}p{2.0cm}>{\raggedright\arraybackslash}p{2.0cm}
  >{\raggedright\arraybackslash}p{2.0cm}>{\raggedright\arraybackslash}p{1.8cm}@{}}
\toprule
test & measured & v0\_4's error & v0\_5's error & role\\
\midrule
energy ledger closure, reference point & & & $\wg{ref.kin.closure}$ of
  \Edep & bookkeeping\\
energy balance at every node of the flow network & & & \wm{nodes} nodes,
  to $10^{-8}$ & bookkeeping\\
\midrule
drift velocity, the \wg{sw.wtr.sets} sets that
  agree~\cite{hunter88,nakamura90,pack62}, \wg{sw.wtr.n} points & &
  \multicolumn{2}{l}{rms \wg{sw.wtr.tr.rms}\,\% (one set of cross
  sections)} & train\\
ionization coefficient, 50 to 300\unit{Td}, the \wg{sw.atr.sets} sets
  that agree~\cite{bhattacharya79,dallarmi92,jacques86,kruithof40,heylen71},
  \wg{sw.atr.n} points & & \multicolumn{2}{l}{rms \wg{sw.atr.tr.rms}\,\%}
  & train\\
drift velocity, the \wg{sw.wh.sets} sets that do not agree (1953 to 1992)
  & & \multicolumn{2}{l}{rms \wg{sw.wh.tr.lo} to \wg{sw.wh.tr.hi}\,\% by
  set} & held out; the two at 195\unit{K} validation\\
characteristic energy $D_T/\mu$, Koizumi et al.~\cite{koizumi86},
  \wg{sw.dtk.n} points, 0.005 to 0.9\unit{Td} & &
  \multicolumn{2}{l}{rms \wg{sw.dtk.tr.rms}\,\%} & validation\\
ionization coefficient below 50\unit{Td}, \wg{sw.alow.sets}
  sets~\cite{specht80,kruithof40,heylen71,dallarmi92,jacques86},
  \wg{sw.alow.n} points & & \multicolumn{2}{l}{mean
  \wg{sw.alow.tr.mean}\,\%} & held out by rule\\
electroluminescence yield, 6 to 13.5\unit{Td}, Mano et
  al.~\cite{mano22} & $\wg{el.mano.a}\,E/N - \wg{el.mano.b}$ &
  \multicolumn{2}{l}{\wg{el.lo} to \wg{el.hi} of it} & held out\\
superelastic return of \lv{5}, \lv{4}, \lv{3} by thermal electrons,
  Blagoev et al.~\cite{blagoev84} & $\wg{blag.1s5.meas}$,
  $\wg{blag.1s4.meas}$, $\wg{blag.1s3.meas}\cm{3}$ & not carried &
  \x{\wg{blag.1s5.lo}}--\wg{blag.1s5.hi}, \x{\wg{blag.1s4.lo}}--\wg{blag.1s4.hi},
  \x{\wg{blag.1s3.lo}}--\wg{blag.1s3.hi} & new, held out\\
ionization coefficient, Monte Carlo on the same cross
  sections~\cite{dias23}, 20 to 50\unit{Td} & & \multicolumn{2}{l}{Biagi's
  file \wg{mc.list}\,\%} & solver check\\
the Monte Carlo's bulk drift velocity and SST coefficient, trained set
  against Biagi's file, 30 to 300\unit{Td} & & \multicolumn{2}{l}{within
  \wg{mc.chk.max}\,\%} & solver check\\
\midrule
energy per ion pair under a beam, from Saito et al.~\cite{saito03} &
  \wg{saito.W}\unit{eV} & \multicolumn{2}{l}{\wg{beam.W.dev}\,\%} &
  held out\\
excitations per ion pair, Saito et al.~\cite{saito03} & \wg{saito.nex} &
  \multicolumn{2}{l}{\wg{beam.nex.dev}\,\%} & held out\\
energy per excitation or ion, Saito et al.~\cite{saito03} &
  \wg{saito.ws}\unit{eV} & \multicolumn{2}{l}{\wg{beam.ws.dev}\,\%} &
  the beam's check\\
VUV efficiency, continuous e-beam, Morozov et al.\ 2008~\cite{morozov08}
  & \wg{moz.meas} $\pm$ \wg{moz.err} & \wg{moz.kin.dev.v04}\,\% &
  \wg{moz.kin.dev}\,\% & held out\\
VUV efficiency at $3\times10^5\unit{W\,cm^{-3}}$, Eckstrom et al.\
  1988~\cite{eckstrom88} & near \wg{eck.meas}, their limit &
  \wg{eck.kin.dev.v04}\,\% & \wg{eck.kin.dev}\,\% at 1\unit{atm} & held
  out\\
\bottomrule
\end{tabularx}
\caption{LIGHT\_Kr v0\_5 against the data, every band at its default, on
the kinetic EEDF, beside v0\_4's errors. v0\_5 runs on v0\_3's and
v0\_4's trained cross sections, so the swarm and the beam's partition
score one set of cross sections. The superelastic row gives v0\_5's rate
over the measurement, for electrons at 300 to 2000\unit{K}. An error is the model less the measurement over the
measurement, signed for a single value and rms (or a mean) over a set.
A validation set is held out of the training but took part in choosing
which factors train (Table~\ref{tab:var}). The swarm rows compare the
model as each experiment measures it: the two-term solution times the
Monte Carlo's ratios (Section~\ref{sec:swarm}). The beam rows run under
the e-beam drive of Section~\ref{sec:beam}.}
\label{tab:score}
\end{table}

\section{Overall accuracy and its uncertainty}
\label{sec:accuracy}

\textbf{Against the data.} Five groups of measurements bear on v0\_5, and
each tests a part of it (Table~\ref{tab:score}).
\begin{itemize}
\item \textbf{The electron swarm} (Section~\ref{sec:swarm}) trains the
cross sections that give every electron rate and scores what the training
did not see. v0\_5 runs on v0\_3's trained set. On the \wg{sw.n.train}
sets that agree with one another the drift velocity is within
\wg{sw.wtr.tr.rms}\,\% rms and the ionization coefficient within
\wg{sw.atr.tr.rms}\,\% from 50 to 300\unit{Td}. Held out, the
characteristic energy of Koizumi et al.\ is within \wg{sw.dtk.tr.rms}\,\%
rms (Biagi's file: \wg{sw.dtk.bi.rms}\,\%), a validation rather than a
test, as it took part in choosing which factors train. The ionization
coefficient below 50\unit{Td}, which nothing in the training saw, lies
\wg{sw.alow.tr.mean}\,\% from the measurements on average, where those
part by up to \x{\wg{sw.low.part}} from one another.
\item \textbf{The electroluminescence yield} (Section~\ref{sec:el}), a
swarm measurement of the light the excitation makes, which no fit saw:
the model lies at \wg{el.lo} to \wg{el.hi} of it from 6 to
13.5\unit{Td}.
\item \textbf{The superelastic return} (Section~\ref{sec:excited}) is a
new score. v0\_5 returns the 5s levels to the ground state by detailed
balance on the trained excitation cross sections; for thermal
electrons that gives \wg{blag.1s5.lo} to \wg{blag.1s5.hi} times Blagoev
et al.'s afterglow rate for \lv{5}, \wg{blag.1s4.lo} to \wg{blag.1s4.hi}
for \lv{4} and \wg{blag.1s3.lo} to \wg{blag.1s3.hi} for \lv{3}.
\item \textbf{The e-beam measurements} (Section~\ref{sec:beam}) score the
chain from excitation to the 147-nm photon. The beam's partition is
v0\_4's: \wg{beam.W.dev}\,\% from Saito et al.\ on the energy per ion pair
and \wg{beam.nex.dev}\,\% on the excitations per pair. At low pump the
error against Morozov et al.\ is \wg{moz.kin.dev}\,\%, \wg{moz.kin.sig}
of their standard errors. At $3\times10^5\unit{W\,cm^{-3}}$ and
1\unit{atm} it is \wg{eck.kin.dev}\,\% against Eckstrom et al.'s report.
\item \textbf{Nothing scores the field-driven regime.} No pure-krypton
measurement in hand tests a gas whose electrons a field heats at the
reference point's density. There v0\_5's answer rests on the electrons'
collisions with excited atoms and with \Krts{}, which v0\_4 held at 0:
their rates are a calculation that no krypton measurement tests, and
their band is propagated below.
\end{itemize}

\textbf{Propagated from the sources.} Each input is set to the two ends
its source allows and the run repeated, the others at their defaults, on
both electron closures (Table~\ref{tab:unc}). The inputs are the
\wg{unc.nband} bands the rows read and the \wg{unc.nrow} measured
constants with a quoted error that no band covers: \Krtp{} + \el{}
($\pm$10\,\%~\cite{shiu77}), \Krp{} + 2Kr (the spread of the three
measurements~\cite{bhattacharya70,grossl81,smith72}) and the \Krts{}
lifetime (270 $\pm$ 5\unit{ns}~\cite{morikawa89}). A calculated row's
band runs from 0 to twice the calculation. Half the difference between
an input's two ends is its term. The trained cross sections are one more
input: the reference point is run at \wg{unc.kin.xs.n} draws of the
training's posterior on the kinetic EEDF and \wg{unc.2t.xs.n} on the
two-temperature closure, and the draws' standard deviation is their
term. The quadrature sum treats the terms as independent standard
deviations, which a band's range is not, so it is indicative. The
excursions from the reference, summed by sign, are the ceiling of the
spread, not a bound. Not propagated: the 5p rung's lines and its levels'
population at the gas temperature (its light is under 1\,\% of the
input), and the excimer's rotational constant, an estimate that scales
its redissociation as $1/r_e^2$ while the well's band moves it
exponentially.

On the 147-nm share at the reference point, \wg{unc.kin.y0}\,\% of the
input on the kinetic EEDF, the quadrature sum is $\pm$\wg{unc.kin.quad}
points ($\pm$\wg{unc.kin.quad.pct}\,\%). The excursions sum to
\wg{unc.kin.down} and \wg{unc.kin.up} points. The largest terms are
\wg{unc.kin.top1} (\wg{unc.kin.top1.t} points), \wg{unc.kin.top2}
(\wg{unc.kin.top2.t}) and \wg{unc.kin.top3} (\wg{unc.kin.top3.t}). The
trained cross sections' term is \wg{unc.kin.xs.sd} points
(\wg{unc.kin.xs.pct}\,\% of the share). On the two-temperature closure
the sum is $\pm$\wg{unc.2t.quad} points ($\pm$\wg{unc.2t.quad.pct}\,\%),
the cross sections' term \wg{unc.2t.xs.sd}.

What dominates has changed again. In v0\_4 it was the $J=1$ levels'
pumping; in v0\_5 the three largest terms are the electrons' collisions
with \Krts, which v0\_4 did not carry, and the first alone makes
\wg{unc.kin.top1.share}\,\% of the quadrature sum's square. The excimer's
ionization, Vriens and Smeets's calculation, takes the share from
\wg{unc.end.ei_kr2.low}\,\% without it to \wg{unc.end.ei_kr2.high}\,\% at
twice the formula. The two rows that nothing measures or calculates move
it as far from their 0: at Kannari et al.'s estimate, $10^{-7}\cm{3}$, the
excimer's quenching by electrons takes it to
\wg{unc.end.eq_kr2.high}\,\%, and the same rate mixing its two states
raises it to \wg{unc.end.em_kr2.high}\,\%, the fast 0$_u^+$ state
radiating in \wm{tauRADF}\unit{ns}, before the electrons can ionize it.
Of the measured inputs, the \lv{4} formation's band, down to a third of
the \lv{5} rate (Audouard et al.'s ratio~\cite{audouard88}), takes the
share to \wg{unc.end.kr2_form_1s4.low}\,\%, and the measured 5s \ra 5p
cross sections' quoted errors span \wg{unc.end.k_5s5p.high} to
\wg{unc.end.k_5s5p.low}\,\%. The $J=1$ levels' proxies, v0\_4's largest
term, now move it the other way and less: with no electron pumping them
the share is \wg{unc.end.k_5s5p_j1.low}\,\%.

\begin{table}[H]
\centering\scriptsize
\resizebox{\linewidth}{!}{%
\begin{tabular}{@{}llllll@{}}
\toprule
input & its two ends & 147\unit{nm}, kinetic & 147\unit{nm}, two-temperature & near IR, kinetic & $T_g$ (K), kinetic\\
\midrule
\wgUncTable
\bottomrule
\end{tabular}}
\caption{The propagated uncertainty's inputs, each alone at its two ends
with the others at their defaults, and the outputs at those ends (low end
to high end), percent of \Edep; for the trained cross sections the mean
and standard deviation over the posterior's draws. With every input at
its default the 147-nm share is \wg{unc.kin.y0}\,\% (kinetic) and
\wg{unc.2t.y0}\,\% (two-temperature).}
\label{tab:unc}
\end{table}

\textbf{Numerical accuracy.}
\begin{itemize}
\item \textbf{The ledger.} The energy ledger closes to
$\wg{ref.kin.closure}$ of the input on the kinetic EEDF and to
$\wg{ref.2t.closure}$ on the two-temperature closure. The energy flows
also balance at each of \wm{nodes} nodes to $10^{-8}$: the electron gas,
each reservoir, each group of rows and each output.
\item \textbf{The kinetic EEDF at the reference point.} The run takes
\wg{ref.kin.n} kinetic solutions, \wg{ref.kin.on} in the pulse, at $E/N$
from \wg{ref.kin.en.lo} to \wg{ref.kin.en.hi}\unit{Td}. Over the pulse's
second half the field carries \wg{ref.kin.ratio.lo} to
\wg{ref.kin.ratio.hi} of the input. There the kinetic and the model's
energy balances agree to $\wg{ref.kin.gap}$ of the input, and the kinetic
mean energy equals the model's to $\wg{ref.kin.reach}$.
\item \textbf{Convergence.} Half the renewal step moves the 147-nm share by
$\wg{conv.kin_step.147}$ of itself. A renewal at least every 5\unit{ns}
moves it by $\wg{conv.kin_max_dt_s.147}$, and twice the energy cells by
$\wg{conv.kin_cells.147}$. A ten times tighter solver tolerance moves it by
$\wg{conv.rtol.147}$, the EEDF renewed \wg{conv.rtol.n} times against
the reference's \wg{ref.kin.n}.
\item \textbf{The afterglow's EEDF.} Through the pulse the kinetic EEDF is
advanced in time from its last solution; after the pulse it takes its
steady shape at each renewal. Advancing it in time after the pulse too
(\wg{ag.n} time steps) moves the 147-nm share by $\wg{ag.147}$ of itself,
the near IR by $\wg{ag.IR}$ and the reaction heat by $\wg{ag.heat}$: the
steady shape costs nothing at this point.
\item \textbf{Restarts.} Where the kinetic EEDF finds no solution from its
warm start, the solver starts it again cold. The mean match has one
solution from any start, so a restart changes no answer, only the way to
it. Over the \wg{kin.runs} kinetic runs of this note the solver restarted
\wg{kin.restarts}. Where a table point told to leave the growth term out
has no steady solution, as in the afterglow of very hot gas, where
Penning ionization makes electrons faster than recombination takes them,
it takes the growth term on the same grid, and failing that a grid of its
own: \wg{kin.fallbacks}. Before v0\_4 such a point stopped the run.
Where even that finds none, v0\_5 leaves the point out of the table
(\wg{kin.dropped}), and where the table's centre has none, LIGHT's
Maxwellian stands in until the next renewal (\wg{kin.failed}); before
v0\_5 either stopped the run, as one did in the afterglow of the scan's
0.25-atm, 500-kW corner.
\item \textbf{The swarm scores} are independent of the grid: half and twice
the \wg{sw.cells} energy cells move the train sets' rms by \wg{sw.grid}
points at most. The emulator the training ran on reproduces the solver to
\wg{emu.W}\,\% in the drift velocity and \wg{emu.a}\,\% in the
ionization coefficient on designs it was not trained on.
\end{itemize}

\section{What v0\_5 is}
\label{sec:what}

LIGHT\_Kr v0\_5 is a 0-D model: one uniform box of pure krypton at
constant volume, into whose electrons a given power is put. It follows
where that energy goes: into the four 5s levels (\Krs: \lv{5}, \lv{4},
\lv{3}, \lv{2}) and the 5p rung (\Krpp), the ions \Krp{} and \Krtp, the
excimer \Krts{} in its slow (1$_u$, 0$_u^-$) and fast (0$_u^+$) states
and its 147-nm second continuum, the 5p \ra 5s lines, the resonance
lines, heat, and the electron gas. Its energy ledger closes by
construction. Every species carries a formation energy, and every row's
energy goes to the electrons, the gas, light or storage.

Its network has \wg{par.rows} rows on \wg{par.species} species
(Appendix~\ref{app:rates}):
\begin{itemize}
\item electron-impact excitation of each of the \wg{par.levels} levels of
Biagi's cross sections at its own threshold, ionization and elastic
loss~\cite{biagi99,pancheshnyi12,carbone21}, with v0\_3's factors trained
on the swarm data (Section~\ref{sec:swarm}): each 5s level makes its own
species, every higher level the 5p rung \Krpp{} (the higher levels
cascading through it, the excess heating the gas);
\item the Hornbeck--Molnar ionization of the \wg{hm.levels} levels at or
above \Krtp's energy~\cite{hornbeck51}, a band at 0;
\item Penning ionization of the 5s levels on one another;
\item ion conversion \Krp{} \ra \Krtp{} and its reverse, and
dissociative recombination of \Krtp{} onto the 5p rung;
\item three-body excimer formation out of each 5s level at its own
rate, its reverse, the 147-nm emission of the two excimer states, and
their quenching by krypton (a band at 0);
\item the transfers between the 5s levels on krypton and their reverses,
and the escaping resonance lines (a band at 0);
\item the rung's relaxation on krypton and its reverse, and its
\wg{par.lines} 5p \ra 5s lines, each from its own level, the rung's ten
levels populated at the gas temperature;
\item the electrons' 5s $\leftrightarrow$ 5p collisions, each of the
\wg{par.pairs} measured pairs~\cite{jung05,jung06} at its own threshold,
and their reverses;
\item the electrons' collisions with the excited atoms and the excimer:
the 5s levels' and the 5p rung's ionization and \Krts's, on Vriens and
Smeets's calculation~\cite{vriens80}; the 5s levels' return to the ground
state; the mixing of \lv{5} and \lv{4}; and, as bands at 0, the
excimer's quenching and the mixing of its two states;
\item and, as bands at 0, the other rows with no measurement.
\end{itemize}
The electrons' rates come from the two-term Boltzmann equation with
electron-electron collisions~\cite{hagelaar05}, solved during the run at
the model's own mean energy (the default, minutes a run), every collision
with an excited population a process of its own in it
(\wg{par.pops} processes). The alternative is a Maxwellian with a
tabulated tail (the two-temperature closure, seconds a run;
Section~\ref{sec:twot}). A rate known only as a published spread, or not
at all, is a band: a default and two ends.

\textbf{What changed from v0\_4.} v0\_4 had \wg{par.rows.v04} rows.
\begin{itemize}
\item \textbf{Calculations where nothing is measured.} The ionization of
the four 5s levels and of the 5p rung by electrons, which v0\_4 held at 0
(no measurement for krypton or argon), is Vriens and Smeets's cross
section~\cite{vriens80} at each level's own ionization energy,
\wg{lad.ion.1s2} to \wg{lad.ion.1s5}\unit{eV} for the 5s levels; the
ionization of \Krts, which v0\_4 did not carry, the same at the
excimer's \wg{lad.ei}\unit{eV}, the order of Flannery et al.'s
binary-encounter calculation~\cite{flannery81}. Each is a band from 0 to
twice the calculation. The excimer's quenching by electrons, e + \Krts{}
\ra{} e + 2Kr, and the mixing of its two states are neither measured nor
calculated in the sources found: bands at 0, to Kannari et al.'s
estimate of the quenching~\cite{kannari85}.
\item \textbf{The pumps' energy.} v0\_4 charged every 5s \ra 5p pumping
the gap to the rung's lowest level, 2p$_{10}$. v0\_5 books each measured
pair as its own row, the electron paying the pair's threshold and the
level's excess over 2p$_{10}$ heating the gas as the rung's levels relax:
on average \wg{jung.gap.1s5}\unit{eV} out of \lv{5} (v0\_4:
\wg{jung.gap0.1s5}), \wg{jung.gap.1s4} out of \lv{4}
(\wg{jung.gap0.1s4}), \wg{jung.gap.1s3} out of \lv{3}
(\wg{jung.gap0.1s3}) and \wg{jung.gap.1s2} out of \lv{2}
(\wg{jung.gap0.1s2}). The rung's ten levels are taken populated at the
gas temperature (an assumption, Section~\ref{sec:limits}); its
superelastic return to each 5s level and its lines are read from each
level, and its relaxation's reverse is detailed balance at the gas
temperature.
\item \textbf{The 5s levels and the electrons.} Each 5s level returns to
the ground state by the Klein--Rosseland reverse of its own excitation
cross section; \lv{5} and \lv{4} mix by electrons at Blagoev and
Pilosov's measured $\wg{blag.pil}\cm{3}$~\cite{blagoev81} (a band to
$\wg{blag.pil.top}$, the electron temperature of their afterglow not
being stated in their abstract).
\item \textbf{The tail table's excited axis.} Section~\ref{sec:twot}.
\item \textbf{Data.} \Krtp{} is bound by Signorell et al.'s
9267.8\unit{cm^{-1}}~\cite{signorell01}; Kholin et al.'s \lv{3}
measurement~\cite{kholin15} has argon for its third body, so \lv{3}'s
excimer formation stays a proxy.
\end{itemize}
The cross sections, the swarm, Jung et al.'s tables and every row v0\_5
did not write are v0\_4's. Table~\ref{tab:undo} undoes each change alone
at the reference point.

\begin{table}[H]
\centering\footnotesize
\begin{tabular}{@{}>{\raggedright\arraybackslash}p{6.0cm}rrrr@{}}
\toprule
undone from v0\_5 & 147\unit{nm}, kinetic & change & two-temperature & $T_g$ end (K)\\
\midrule
none (v0\_5) & \wg{ref.kin.147} & & \wg{ref.2t.147} & \wg{ref.kin.tg}\\
\wgUndo
\midrule
v0\_4 & \wg{chg.kin.v04} & & \wg{chg.2t.v04} & \wg{chg.kin.tg.v04}\\
\bottomrule
\end{tabular}
\caption{What each of v0\_5's changes does at the reference point, and
two of v0\_4's: each undone alone, the others kept, percent of \Edep;
change is relative to v0\_5 on the kinetic EEDF. v0\_4 is the last row.
The pumps' energy cannot be undone alone; the text gives its size.}
\label{tab:undo}
\end{table}

The calculated rows are most of the change. With the three at v0\_4's 0
the share is \wg{undo.no_calc.kin}\,\%, below v0\_4's
\wg{chg.kin.v04}\,\%: the rest of v0\_5 lowers it, chiefly the changes
that cannot be undone alone, the pumps' energy, each pair charged its own
threshold (\wg{ref.kin.pumpx}\,\% of the input to heat), and the rung's
levels at the gas temperature. The excimer's
ionization alone moves the share by \wg{undo.no_ei.kin.back}\,\%:
without it the share is \wg{undo.no_ei.kin}\,\%. The stepwise ionizations
move it by \wg{undo.no_kstep.kin.back}\,\%, the superelastic return by
\wg{undo.no_super.kin.back}\,\%, the electrons' 1s$_5$ \ra 1s$_4$ by
\wg{undo.no_emix.kin.back}\,\%, and the rung's reverse by
\wg{undo.no_kpsr.kin.back}\,\% at the reference point's
\wg{ref.kin.tg}\unit{K}. With every level forming the excimer at Ueda's
rate, v0\_3's, the share is \wg{undo.ueda.kin}\,\%, and without the
excimer's redissociation \wg{undo.no_redissoc.kin}\,\%. The pumping of the
two $J=1$ levels now raises the share (\wg{undo.no_j1.kin}\,\% without
it), where in v0\_4 it lowered it.

\section{The reference point}
\label{sec:ref}

LIGHT's drive in pure krypton: \wm{p}\unit{atm} at \wm{Tg0}\unit{K},
\wm{P}\unit{kW\,cm^{-3}} for \wm{pulse}\unit{\mu s} into the electrons, a
$\wm{ne}\unit{cm^{-3}}$ seed, \wm{tend}\unit{\mu s} in all.
Table~\ref{tab:ref} gives its outputs on both closures.

\begin{table}[H]
\centering\footnotesize
\begin{tabular}{@{}>{\raggedright\arraybackslash}p{8.6cm}rr@{}}
\toprule
 & kinetic EEDF (default) & two-temperature\\
\midrule
147\unit{nm}, \Krts{} & \wg{ref.kin.147} & \wg{ref.2t.147}\\
near IR, 5p \ra 5s & \wg{ref.kin.IR} & \wg{ref.2t.IR}\\
heat, reactions and relaxation & \wg{ref.kin.heat} & \wg{ref.2t.heat}\\
heat, elastic collisions & \wg{ref.kin.elastic} & \wg{ref.2t.elastic}\\
stored in ions (the seed's ions recombine) & \wg{ref.kin.ions} & \wg{ref.2t.ions}\\
to \Krs{} and \Krpp{} by electron impact, gross & \wg{ref.kin.krstar} & \wg{ref.2t.krstar}\\
to \Krp{} & \wg{ref.kin.krplus} & \wg{ref.2t.krplus}\\
to \Krts{} (formed \x{} \wm{hv.RAD}\unit{eV}) & \wg{ref.kin.kr2} & \wg{ref.2t.kr2}\\
the 5s \ra 5p pumps' excess over 2p$_{10}$, to heat (gross) & \wg{ref.kin.pumpx} & \wg{ref.2t.pumpx}\\
\Krs{} formed that becomes \Krts{} (\%) & \wg{ref.kin.ak8} & \wg{ref.2t.ak8}\\
\Krs{} formed that the electrons ionize (\%) & \wg{ref.kin.krs.step} & \wg{ref.2t.krs.step}\\
\Krs{} formed that returns to the ground state by electrons (\%) & \wg{ref.kin.krs.super} & \wg{ref.2t.krs.super}\\
\Krpp{} formed that the electrons ionize (\%) & \wg{ref.kin.kpp.step} & \wg{ref.2t.kpp.step}\\
\Krts{} formed out of \lv{5}, \lv{4}, \lv{3} (\%) & \wg{ref.kin.form.1s5}, \wg{ref.kin.form.1s4}, \wg{ref.kin.form.1s3} & \wg{ref.2t.form.1s5}, \wg{ref.2t.form.1s4}, \wg{ref.2t.form.1s3}\\
\Krts{} formed that radiates (\%) & \wg{ref.kin.k2rad} & \wg{ref.2t.k2rad}\\
\Krts{} formed that the electrons ionize (\%) & \wg{ref.kin.k2ei} & \wg{ref.2t.k2ei}\\
\Krts{} formed that redissociates (\%) & \wg{ref.kin.k2redis} & \wg{ref.2t.k2redis}\\
ions made by \Krts's ionization (\% of all, by number) & \wg{ref.kin.ion.ei} & \wg{ref.2t.ion.ei}\\
\quad by the 5s levels' ionization & \wg{ref.kin.ion.step_s} & \wg{ref.2t.ion.step_s}\\
\quad by the 5p rung's ionization & \wg{ref.kin.ion.step_p} & \wg{ref.2t.ion.step_p}\\
\quad by Penning ionization & \wg{ref.kin.ion.pen} & \wg{ref.2t.ion.pen}\\
\quad from the ground state & \wg{ref.kin.ion.C2} & \wg{ref.2t.ion.C2}\\
5s \ra 5p pumpings per \Krts{} formed, net & \wg{ref.kin.cyc} & \wg{ref.2t.cyc}\\
$T_e$, pulse mean (eV) & \wg{ref.kin.te} & \wg{ref.2t.te}\\
$n_e$, pulse mean (cm$^{-3}$) & $\wg{ref.kin.ne.on}$ & $\wg{ref.2t.ne.on}$\\
$T_g$ at the end (K) & \wg{ref.kin.tg} & \wg{ref.2t.tg}\\
$n_e$ at the end (cm$^{-3}$) & $\wg{ref.kin.ne}$ & $\wg{ref.2t.ne}$\\
\bottomrule
\end{tabular}
\caption{The reference point at $t=2\unit{\mu s}$, percent of \Edep{}
unless a row says otherwise. The fate rows (the first five, with the
other stores) sum to 100; the flows below them overlap.}
\label{tab:ref}
\end{table}

The electrons pump \Krs{} to the 5p rung \wg{ref.kin.cyc} times for each
\Krts{} formed net, and krypton relaxes the rung back, so the cycle's
thresholds, charged pair by pair, become heat: \wg{ref.kin.cycheat}\,\%
of the input, \wg{ref.kin.perphoton}\unit{eV} for each 147-nm photon. Of
the pumping, \wg{ref.kin.pump.1s5}\,\% is out of \lv{5},
\wg{ref.kin.pump.1s4}\,\% out of \lv{4}, \wg{ref.kin.pump.1s2}\,\% out of
\lv{2} and \wg{ref.kin.pump.1s3}\,\% out of \lv{3}. The excimers that form
mostly do not radiate: the electrons ionize \wg{ref.kin.k2ei}\,\% of them,
\wg{ref.kin.k2rad}\,\% radiate and \wg{ref.kin.k2redis}\,\% redissociate.
Counted by number, the excimer's ionization makes \wg{ref.kin.ion.ei}\,\%
of the ions, the 5s levels' \wg{ref.kin.ion.step_s}\,\%, the rung's
\wg{ref.kin.ion.step_p}\,\%, Penning ionization \wg{ref.kin.ion.pen}\,\%
and the ground state \wg{ref.kin.ion.C2}\,\%, and every ion recombines
onto the 5p rung (R11p): an excimer ionized is an atom sent round the
cycle again. A stepwise ionization costs the electron a fraction of the
ground state's \wm{E.ion}\unit{eV}, and the pulse holds
$\wg{ref.kin.ne.on}\unit{cm^{-3}}$ of electrons (v0\_4:
$\wg{chg.kin.ne.v04}$) at \wg{ref.kin.te}\unit{eV} (v0\_4:
\wg{chg.kin.te.v04}). On the kinetic EEDF the share differs from the
two-temperature closure's by \wg{ref.kin.vs2t}\,\%. The cycle's measured
cross sections, scaled within their errors and against Kannari's fit in
their place~\cite{kannari85,jacob76}:

\begin{table}[H]
\centering\footnotesize
\begin{tabular}{@{}lrrrrr@{}}
\toprule
5s \ra 5p cross sections & 0 (no cycle) & \x{} \wm{bd.k_5s5p.lo} & \x{} 1, Jung et al. & \x{} \wm{bd.k_5s5p.hi} & Kannari's fit\\
\midrule
147\unit{nm}, kinetic EEDF & \wg{k5.0.kin} & \wg{k5.0.6.kin} & \wg{k5.1.kin} & \wg{k5.1.4.kin} & \wg{kann.kin}\\
147\unit{nm}, two-temperature & \wg{k5.0.2t} & \wg{k5.0.6.2t} & \wg{k5.1.2t} & \wg{k5.1.4.2t} & \wg{kann.2t}\\
\bottomrule
\end{tabular}
\caption{The 147-nm share (percent of \Edep) against the scale on the
measured 5s \ra 5p cross sections, every pair's. The band runs from
\wm{bd.k_5s5p.lo} to \wm{bd.k_5s5p.hi}, the measurements' quoted errors;
the two $J=1$ levels' proxies have a band of their own
(\texttt{k\_5s5p\_j1}, Table~\ref{tab:unc}).}
\label{tab:k5}
\end{table}

Jung et al.\ measured the excitation out of the two metastable levels,
\lv{5} and \lv{3}, into the ten 5p levels, with stated uncertainties of
$\pm$\wg{jung.unc.1s5} and $\pm$\wg{jung.unc.1s3}\,\%~\cite{jung06}. For
the two $J=1$ levels nothing is measured; each takes the set of the
metastable with the same ion core, shifted to its own threshold (Zeng et
al.'s R-matrix calculation covers them~\cite{zeng05}; its numbers are not
in a source in hand). At the reference point's \wg{jung.te.ref}\unit{eV}
the four levels' Maxwellian rates, each the sum of its pairs, are
$\wg{jung.k.1s5}$, $\wg{jung.k.1s4}$, $\wg{jung.k.1s3}$ and
$\wg{jung.k.1s2}\cm{3}$ (\lv{5}, \lv{4}, \lv{3}, \lv{2}); the pairs' sum
is v0\_4's level rate to $\wg{jung.sum.dev}$ of it (their onsets at Biagi's
level energies, where Table~I of the measurement prints two decimals).

Pressure and power move the share as much (Table~\ref{tab:scan},
two-temperature closure). It rises with pressure and falls with power at
every point of the scan, from nothing at 0.1\unit{atm} to
\wg{scan.best}\,\% at \wg{scan.best.atm}\unit{atm} and
\wg{scan.best.kw}\unit{kW\,cm^{-3}}. At 0.1\unit{atm} a microsecond of
this drive heats the gas to \wg{scan.tg01.lo} to
\wg{scan.tg01.hi}\unit{K}, beyond the measured range of several rates,
and at 500\unit{kW\,cm^{-3}} \wg{scan.ions01}\,\% of the input is still
in ions at the end.

\begin{table}[H]
\centering\footnotesize
\begin{tabular}{@{}lrrrrr@{}}
\toprule
kW\,cm$^{-3}$ & 0.1\unit{atm} & 0.25\unit{atm} & 0.5\unit{atm} & 1\unit{atm} & 2\unit{atm}\\
\midrule
\wgScan
\bottomrule
\end{tabular}
\caption{The 147-nm share (percent of \Edep) over pressure and power,
1\unit{\mu s} pulses, the two-temperature closure on its field-driven
tail table.}
\label{tab:scan}
\end{table}

\section{The two-temperature closure and its tail}
\label{sec:twot}

The two-temperature closure keeps the electrons Maxwellian at $T_e$ below
krypton's lowest threshold, \wm{E.S}\unit{eV}, and lets the distribution
fall above it with a tail temperature $T_t = r\,T_e$, continuous at the
break. It needs $r$, and takes it from a table built from the model's own
kinetic EEDF for a field-driven discharge: at each point, ground-state
krypton at 300\unit{K} and 1\unit{atm}, as many ions as electrons, a field
at the mean energy $\frac32 T_e$, electron-electron collisions at $n_e/N$
and the growth term, as the run solves them; $r$ is the tail ratio at
which the two-temperature distribution removes the same energy through
the 51 levels and the ionization as the kinetic one.

\textbf{The excited axis.} v0\_4's table, $r(n_e/N, T_e)$, had no excited
atoms in its states, and along v0\_4's runs its tail lay below the
kinetic EEDF's own. v0\_5's is $r(n_e/N, T_e, x_s)$ on \wg{tail.npts}
points: $\log_{10} n_e/N$ from $\wg{tail.x.lo}$ to $\wg{tail.x.hi}$,
$T_e$ from \wg{tail.te.lo} to \wg{tail.te.hi}\unit{eV}, and the 5s
fraction $x_s=[\Krs]/[\mathrm{Kr}]$ at \wg{tail.nxs} values from 0 to
$\wg{tail.xs.top}$. At each point the 5s levels hold $x_s$ by their
weights, and the 5p rung and \Krts{} sit in their steady state against
them at the point's $T_e$ and $n_e$, from the model's own rates. Its
$x_s=0$ block is v0\_4's table to the printed digit. The excited atoms'
superelastic electrons heat the tail: at $n_e/N=10^{-5}$ and
2.5\unit{eV}, $r$ is \wg{tail.ref.r0} without them and \wg{tail.ref.r3}
at $x_s=10^{-3}$. Where the kinetic tail is colder than the form can
follow (\wg{tail.nfloor} points) $r$ takes the floor, $10^{-3}$; where it
is hotter (\wg{tail.ntop} points, at low $T_e$ with excited atoms, where
the superelastic electrons put far more than the Maxwellian's loss into
the tail), $r$ takes the top, 3, the closest the form comes. Half and
twice the pressure move $r$ by at most \wg{tail.p.dev}\,\% at the points
the runs visit, and 700\unit{K} gas by \wg{tail.tg.dev}\,\%.

Along the reference run the table at the run's 5s fraction
($\wg{tail.xs.lo}$ to $\wg{tail.xs.hi}$) gives \wg{tail.f.lo} to
\wg{tail.f.hi} of the kinetic EEDF's own tail ratio; its $x_s=0$ block,
v0\_4's table, gives \wg{tail.f0.lo} to \wg{tail.f0.hi}, and v0\_3's beam
table \wg{tail.b.lo} to \wg{tail.b.hi} (Table~\ref{tab:tail}). Over the
scan's corners the table's ratio is \wg{tail.cr.f.lo} to
\wg{tail.cr.f.hi} of the kinetic one. The closure's 147-nm share then
lies \wg{tail.ref.f.dev}\,\% from the kinetic EEDF's at the reference
point, where on the $x_s=0$ block it lay \wg{tail.ref.f0.dev}\,\%: with
the tail right, what is left is the Maxwellian below the break, where the
stepwise ionizations' thresholds lie, and the zero block's hotter tail
had more than offset it. At the scan's corners the
closure is within \wg{tail.worst.f}\,\% of the kinetic EEDF, and v0\_3's
beam table, made for an electron beam, is \wg{tail.ref.b.dev}\,\% off at
the reference point.

\begin{table}[H]
\centering\footnotesize
\resizebox{\linewidth}{!}{%
\begin{tabular}{@{}rrrrrrrrr@{}}
\toprule
 & & \multicolumn{3}{c}{147\unit{nm} (\% of \Edep)} & \multicolumn{3}{c}{$r$ over the kinetic's} & \\
\cmidrule(lr){3-5}\cmidrule(lr){6-8}
atm & kW\,cm$^{-3}$ & kinetic & field table & v0\_3's table & field & its $x_s=0$ & v0\_3's & $T_g$ end (K)\\
\midrule
\wgTail
\bottomrule
\end{tabular}}
\caption{The two-temperature closure against the kinetic EEDF at the
scan's corners and the reference point: the 147-nm share on the kinetic
EEDF and on the closure with the field table and with v0\_3's beam table
(the deviation in brackets), and the tables' tail ratio over the kinetic
run's own, averaged over the pulse body (50\unit{ns} to the pulse's end):
the field table at the run's $x_s$, its $x_s=0$ block, and v0\_3's.}
\label{tab:tail}
\end{table}

\section{The electrons and the excited atoms}
\label{sec:excited}

v0\_5 lets the electrons collide with every excited population: the four
5s levels, the ten levels of the 5p rung and both states of \Krts. Under
the kinetic EEDF each collision is a process of its own
(\wg{par.pops} in all), and on a solved EEDF each gives the electrons, or
takes from them, its rate times its energy, a check every process passes.

\textbf{Ionization of the excited atoms and of \Krts.} No measurement of
it for krypton or argon was found. v0\_5 takes Vriens and Smeets's
semi-empirical cross section~\cite{vriens80},
\[\sigma(E)=\frac{4\pi a_0^2R^2}{E+3.25\,E_n}\left(\frac{5}{3E_n}
-\frac1E-\frac{2E_n}{3E^2}\right),\]
at each level's own ionization energy $E_n$: \wg{vs.1s5.En}\unit{eV} for
\lv{5} down to \wg{lad.ion.1s2}\unit{eV} for \lv{2}, \wg{vs.2p10.En}\unit{eV}
for 2p$_{10}$, and \wg{vs.kr2.En}\unit{eV} for \Krts, which it ionizes to
\Krtp. Over a Maxwellian it gives their rate formula to within
$\wg{vs.int.dev}$ for $T_e$ of 1 to 5\unit{eV} (\wg{vs.int.dev05}\,\% at
0.5\unit{eV}). Its peak is $\wg{vs.1s5.peak}\unit{cm^2}$ at
\wg{vs.1s5.peakE}\unit{eV} for \lv{5}, $\wg{vs.2p10.peak}$ at
\wg{vs.2p10.peakE}\unit{eV} for 2p$_{10}$ and $\wg{vs.kr2.peak}$ at
\wg{vs.kr2.peakE}\unit{eV} for \Krts. Flannery et al.'s binary-encounter
calculation gives \Krts{} about $10^{-15}\unit{cm^2}$, 40 to 50\,\% above
the metastable atom's~\cite{flannery81}; the formula puts it
\wg{vs.kr2.over}\,\% above \lv{5}'s. Kannari et al.\ print fits of the
same three rates~\cite{kannari85}, from Vriens's earlier
calculation~\cite{vriens65} and from Vriens and Smeets's: at the
reference point's \wg{jung.te.ref}\unit{eV} they are \wg{kann.s.ref},
\wg{kann.p.ref} and \wg{kann.kr2.ref} times the formula (5s, 5p,
\Krts), and \wg{kann.lo} to \wg{kann.hi} times it from 1 to
3\unit{eV}. Each row's band runs from 0 to twice the formula.

At the reference point the electrons ionize \wg{ref.kin.krs.step}\,\% of
the 5s atoms formed, \wg{ref.kin.kpp.step}\,\% of the rung's and
\wg{ref.kin.k2ei}\,\% of the excimers formed (\wg{ref.kin.k2rad}\,\%
radiate). In the pulse, on a Maxwellian at its mean $T_e$ and $n_e$, the
formula ionizes an excimer in about \wg{ei.tau}\unit{ns}, against its
\wm{tauRAD}-ns radiative lifetime.

\textbf{The superelastic return.} Each 5s level returns to the ground
state by the Klein--Rosseland reverse of its own excitation cross section,
the electron gaining the level's energy. Blagoev et al.\ measured these
rates for thermal electrons in a krypton afterglow~\cite{blagoev84}:
$\wg{blag.1s5.meas}$, $\wg{blag.1s4.meas}$ and
$\wg{blag.1s3.meas}\cm{3}$ for \lv{5}, \lv{4} and \lv{3}. For Maxwellian
electrons at 300 to 2000\unit{K} the model's rates are \wg{blag.1s5.lo}
to \wg{blag.1s5.hi}, \wg{blag.1s4.lo} to \wg{blag.1s4.hi} and
\wg{blag.1s3.lo} to \wg{blag.1s3.hi} times theirs (at 300\unit{K},
$\wg{blag.1s5.k300}$, $\wg{blag.1s4.k300}$ and
$\wg{blag.1s3.k300}\cm{3}$). Such electrons read the excitation cross
sections within a few tenths of an eV of their thresholds, where their
shape is Biagi's: the training scales each group of cross sections as a
whole. The measurement scores the returns and sets nothing. At the
reference point they return \wg{ref.kin.krs.super}\,\% of the 5s atoms
formed and move the share by \wg{undo.no_super.kin.back}\,\%.

\textbf{The 5s levels mixed by electrons.} Blagoev and Pilosov measured
\el{} + \lv{5} \ra \el{} + \lv{4} at $\wg{blag.pil}\cm{3}$ in a krypton
afterglow~\cite{blagoev81}, the electron temperature not given in their
abstract. v0\_5 takes it as $A\,e^{-E_\mathrm{g}/T_e}$ over the levels'
gap $E_\mathrm{g}=\wg{lad.mix15}\unit{eV}$, with $A$ their value (the band's low
end, as if their electrons were hot against the gap) up to
$\wg{blag.pil.top}$ (their value taken at 0.07\unit{eV}, the coldest of
Qiu et al.'s afterglow~\cite{qiu18}), the reverse by detailed balance. At
the reference point's mean $n_e$ and $T_e$ it turns \lv{5} into \lv{4}
\wg{emix.vs.form} times as fast as \lv{5} forms the excimer, and removing
it moves the share by \wg{undo.no_emix.kin.dev}\,\%.

\textbf{The rows at 0.} \el{} + \Krts{} \ra \el{} + 2Kr and the mixing of
the excimer's two states by electrons are neither measured nor calculated
in the sources found. Their bands run from 0 to Kannari et al.'s estimate
of the quenching, $10^{-7}\cm{3}$~\cite{kannari85}; under the kinetic EEDF
each enters as a pseudo cross section whose Maxwellian rate is its rate
law exactly.

\textbf{The 5p rung's levels.} The rung's ten levels are populated at the
gas temperature: at 300\unit{K} \wg{lad.rung.f10.300}\,\% of it is
2p$_{10}$ (partition function \wg{lad.rung.z.300}), at 1000\unit{K}
\wg{lad.rung.f10.1000}\,\%. A pump to a level above 2p$_{10}$ pays its own
threshold; the excess heats the gas as the levels relax, and the gas
lends it back to a reverse, an ionization or a line out of that level. At
300\unit{K} the rung's lines give it a lifetime of
\wg{lad.rung.tau300}\unit{ns}, nearly 2p$_{10}$'s alone (v0\_4:
\wg{tau5p}\unit{ns}, the levels' mean by weight). Its relaxation on
krypton has its detailed-balance reverse, which reaches
\wg{kpsr.max}\,\% of it over the scan's run at \wg{kpsr.max.atm}\unit{atm}
and \wg{kpsr.max.kw}\unit{kW\,cm^{-3}}, and $\wg{kpsr.ref}$ of it at the
reference point's conditions (two-temperature closure).

\section{The electron swarm, and what the training did}
\label{sec:swarm}

The kinetic EEDF's solver is run as a swarm: no electron-electron
collisions, the temporal growth term on, and every grid sized so that the
EEDF falls 11 to 25 decades across it. It gives the drift velocity (the
Joule power over $eE$), the ionization coefficient $\alpha/N=k_\mathrm{i}/W$
and the characteristic energy $D_T/\mu$, with the transverse diffusion
taken over the same momentum transfer and growth term as the field's
flux~\cite{hagelaar05}. Swarm measurements are measurements of the cross
sections themselves, so they may train them, while the device and beam
measurements of Table~\ref{tab:score} may not. v0\_3 trained a few
factors on Biagi's level-resolved cross sections against the swarm data
that agree with one another, and v0\_4 and v0\_5 run on them unchanged;
this section is v0\_3's, its numbers computed again on v0\_5's code.

\textbf{The data.} \wg{sw.nsets} datasets, \wg{sw.npts} points, from
Dutton's survey~\cite{dutton75} and LXCat's LAPLACE
database~\cite{pancheshnyi12,carbone21}, which digitized them from the
papers' tables or figures (Table~\ref{tab:swdata}). None prints an
error, so each carries a class value: 3\,\% for a drift velocity measured
after 1985, 5\,\% before, 7\,\% for a characteristic energy and 10\,\% for
an ionization coefficient. A 2\,\% error in $E/N$ enters through the
model's own slope, which matters where the drift velocity climbs out of
the Ramsauer minimum (0.01 to 0.03\unit{Td}) and for the ionization
coefficient at low field. A dataset's $\chi^2$ per point near 1 means the
model agrees with it within these errors. Points above 300\unit{Td} are
checks only: there the mean energy passes 7\unit{eV}, where the two-term
approximation and the growth model part ways. The ionization coefficient
below 50\unit{Td} never trains (below).

\begin{table}[H]
\centering\scriptsize
\renewcommand{\arraystretch}{1.08}
\begin{tabular}{@{}llrrlrrrrr@{}}
\toprule
 & & $E/N$ & & & error & \multicolumn{2}{c}{$\chi^2$ per point} & \multicolumn{2}{c}{rms (\%)}\\
dataset & quantity & (Td) & points & role & (\%) & Biagi & trained & Biagi & trained\\
\midrule
\wgSwTable
\bottomrule
\end{tabular}
\caption{Krypton's electron swarm data, their roles and the model's
error on Biagi's file as published and on the trained set (v0\_3's to
v0\_5's), each as the experiment measures it. Drift velocities $W$ at
300\unit{K} unless labelled; Hunter et al.'s at 301\unit{K}. A
validation set is held out of the training but took part in choosing
which factors train (Table~\ref{tab:var}).}
\label{tab:swdata}
\end{table}

\textbf{What a swarm experiment measures.} The two-term solution with the
temporal growth term gives the flux drift velocity and $k_\mathrm{i}/W$. A
drift tube times the swarm's centre of mass, the bulk drift velocity, and
a steady-state Townsend experiment measures the spatial growth of its
current, the SST coefficient. A Monte Carlo solution of the Boltzmann
equation, with no two-term approximation, gives all of
them~\cite{dias23}. Run on Biagi's file as published
(Table~\ref{tab:mc}), it finds the bulk drift velocity
\wg{mc.w300}\,\% above the flux one at 300\unit{Td}, and the SST
coefficient \wg{mc.a300}\,\% below the two-term's there. The two-term flux
drift velocity matches the Monte Carlo's to \wg{mc.flux.w}\,\% from 10 to
300\unit{Td}, and from 40\unit{Td} up the two-term $k_\mathrm{i}/W$ lies
\wg{mc.flux.a.lo} to \wg{mc.flux.a.hi}\,\% below the Monte Carlo's flux
coefficient. Below 20\unit{Td}, without ionization, the bulk and flux
drift velocities are one, and the two-term drift velocity is compared as
it is; its own error there (\wg{mc.flux.wlow}\,\% at
\wg{mc.flux.wlow.en}\unit{Td}) is absorbed by the trained factors, which
are effective for the two-term solver. The two-term transverse
diffusion is
\wg{mc.dt.lo} to \wg{mc.dt.hi}\,\% too high from 10 to 300\unit{Td},
so $D_T/\mu$ is scored only below 1\unit{Td}, and Al-Amin and Lucas's
characteristic energies, 1.4\unit{Td} and up~\cite{alamin87}, are a check
of the correction rather than of the cross sections. The model's two-term
solution times these ratios is compared with each measurement. The ratios
are taken on Biagi's file; on the trained set the Monte Carlo gives the
drift velocity's and the ionization coefficient's to \wg{mc.chk.max}\,\%.
Without them the model would be scored, and trained, against the wrong
quantity: comparing the two-term coefficient with Townsend experiments
puts it closer to them at high field than it is.

\begin{table}[H]
\centering\footnotesize
\begin{tabular}{@{}rrrrrr@{}}
\toprule
 & \multicolumn{3}{c}{Biagi's file} & \multicolumn{2}{c}{the trained set}\\
$E/N$ (Td) & $W_\mathrm{bulk}$ & $\alpha_\mathrm{SST}$ & $D_T/\mu$ & $W_\mathrm{bulk}$ & $\alpha_\mathrm{SST}$\\
\midrule
\wgMcTable
\bottomrule
\end{tabular}
\caption{The Monte Carlo's value over the two-term solution's, LoKI-MC at
300\unit{K} and 1\unit{Torr} with $10^5$ electrons~\cite{dias23}. The
coefficient below 40\unit{Td} takes the 40-Td ratio, where the Monte
Carlo counts too few ionizations to resolve it.}
\label{tab:mc}
\end{table}

\textbf{The training.} The factors multiply Biagi's cross sections: the
elastic momentum transfer by a factor read log-linearly between six
energies (0.03 to 20\unit{eV}, held beyond them), and the excitation of
the four 5s levels, the ten 5p levels and the 37 higher levels by one
factor each. The factors at \wg{fac.held}\unit{eV} stay at 1 (below),
so \wg{fac.n} train. The ionization cross section is not trained: beam
experiments measure it, and in a swarm its scale trades against the
excitation's. The prior on each factor's logarithm is broad (standard
deviation 0.5, a factor 1.65), so the data decide. Each trial of the
factors needs the solver at every measured point, so the fits run on a
neural network trained to reproduce it: \wg{emu.n} sets of factors drawn
over \x{\wg{fac.box}} to \x{\wg{fac.box.hi}} (a Latin
hypercube~\cite{mckay79}), each solved at \wg{emu.nodes} values of $E/N$
and gas temperature, and a three-layer network fitted to the change of
each quantity from Biagi's file (Adam~\cite{kingma14}). On designs it was
not trained on it reproduces the solver to \wg{emu.W}\,\% in the drift
velocity, \wg{emu.DT}\,\% in $D_T/\mu$ and \wg{emu.a}\,\% in the
ionization coefficient, and every answer below is re-checked on the
solver itself.

\emph{Which data train.} The consistency check fits every candidate
together, scores each against that joint fit, and drops the worst while
its $\chi^2$ per point exceeds 3. Of \wg{cons.cands} candidates it drops
\wg{cons.ndropped}: \wg{cons.dropped} (their $\chi^2$ when dropped). The
\wg{sw.n.train} that remain train. The choice of 3 matters little below
it: at 2 the same sets remain. At 4, \wg{cons.at4} join, and the 1960 drift
velocities, \wg{bowe.hunter}\,\% from the 1988 ones where both exist on
average, pull the momentum
transfer at 6\unit{eV} to \x{\wg{cons.at4.m6}} and the 5s excitation to
\x{\wg{cons.at4.x1s}}.

\emph{Which factors train.} Two of the nine stay at 1, by a choice made
on held-out data (Table~\ref{tab:var}). With all nine free the fit
matches the train sets best ($\chi^2$ \wg{var.all.train} per point
pooled, against \wg{var.v3.train}), but it lowers the momentum transfer
at the Ramsauer minimum to \x{\wg{var.all.m06}} and raises it at
0.03\unit{eV} to \x{\wg{var.all.m003}}, a trade the train sets allow and
the data that weigh those energies differently reject: Koizumi et al.'s
characteristic energy~\cite{koizumi86} goes from $\chi^2$
\wg{var.bi.dtk} per point on Biagi's file to \wg{var.all.dtk}, and Pack
et al.'s drift velocities in gas at 195\unit{K}~\cite{pack62,pack92} from
\wg{var.bi.p62} to \wg{var.all.p62} (1962) and from \wg{var.bi.p92} to
\wg{var.all.p92} (1992). With the
factors at \wg{fac.held}\unit{eV} held at 1, every train set stays
within the consistency threshold (Hunter et al.'s drift velocity at
\wg{var.v3.hunter} rather than \wg{var.all.hunter}), and those three fit
better than on Biagi's file, at \wg{var.v3.dtk}, \wg{var.v3.p62} and
\wg{var.v3.p92}, the best of the five variants. Holding the 0.6-eV
factor too, or every momentum-transfer factor, fits both the train sets
and these three worse. v0\_3 trained that variant's \wg{fac.n} factors.
Because these held-out sets chose it, they validate the training rather
than test it. The tests that neither the training nor this choice saw
are the ionization coefficient below 50\unit{Td} and the beam's
partition (Section~\ref{sec:beam}).

\begin{table}[H]
\centering\footnotesize
\begin{tabular}{@{}lrrrrrr@{}}
\toprule
 & \multicolumn{3}{c}{train sets} & \multicolumn{3}{c}{held out (validation)}\\
\cmidrule(lr){2-4}\cmidrule(l){5-7}
 & & & Hunter & Koizumi & Pack 1962, & Pack 1992,\\
factors trained & pooled & worst & 1988 & 1986 & 195\unit{K} & 195\unit{K}\\
\midrule
\wgVarTable
\bottomrule
\end{tabular}
\caption{Which factors train: each variant's least-squares factors on
the train set, and the $\chi^2$ per point on the solver itself, each
dataset as its experiment measures it with Table~\ref{tab:swdata}'s
errors. Pooled is over the train sets' \wg{var.npts} points. v0\_3 to
v0\_5 run on the fourth row's factors.}
\label{tab:var}
\end{table}

\emph{The fit.} Least squares on the train set gives the factors; an
ensemble MCMC~\cite{goodman10} (acceptance \wg{mcmc.acc}) gives their
posterior, and a neural posterior estimator, a network trained on
simulated datasets to read the factors off them~\cite{stokes20}, gives
the same answer independently (Table~\ref{tab:fac}). Its $\pm1\sigma$
intervals cover \wg{nn.cov.lo} to \wg{nn.cov.hi} of held-out simulations
(ideal 0.68). The data pin the momentum transfer from
\wg{fac.free.lo} to \wg{fac.free.hi}\unit{eV} and the 5s and higher
excitations (the network's $R^2$ at least \wg{nn.r2.next}); they barely
pin the 5p excitation ($R^2$ \wg{nn.r2.lo}).

\begin{table}[H]
\centering\footnotesize
\begin{tabular}{@{}lrrlrr@{}}
\toprule
factor on & least squares & MCMC median & 16--84\,\% & network & $R^2$\\
\midrule
\wgFacTable
\bottomrule
\end{tabular}
\caption{The trained factors on Biagi's cross sections. v0\_3 to v0\_5
run on the least-squares values. The network column gives its mean and the standard
deviation of the factor's logarithm; $R^2$ is the network's on held-out
simulations.}
\label{tab:fac}
\end{table}

\begin{figure}[H]
\centering
\begin{tikzpicture}
\begin{axis}[width=5.5cm, height=5.6cm, xmode=log, ymode=log,
  xlabel={$E/N$ (Td)}, ylabel={$W$ (cm\,s$^{-1}$)},
  tick label style={font=\scriptsize}, label style={font=\footnotesize},
  legend style={font=\tiny, draw=none, fill=none, at={(0.03,0.97)},
  anchor=north west}, legend cell align=left]
\plotcoords{cE, thick}{\wgSwWtr}
\addlegendentry{trained}
\plotcoords{ink2, dashed}{\wgSwWbi}
\addlegendentry{Biagi's file}
\plotcoords{only marks, mark=*, mark size=1.1pt, ink}{\wgSwWTdata}
\addlegendentry{train}
\plotcoords{only marks, mark=o, mark size=1.1pt, ink2}{\wgSwWHdata}
\addlegendentry{held out}
\end{axis}
\end{tikzpicture}\hfill
\begin{tikzpicture}
\begin{axis}[width=5.5cm, height=5.6cm, xmode=log, ymode=log,
  xlabel={$E/N$ (Td)}, ylabel={$\alpha/N$ (cm$^2$)},
  tick label style={font=\scriptsize}, label style={font=\footnotesize},
  legend style={font=\tiny, draw=none, fill=none, at={(0.97,0.03)},
  anchor=south east}, legend cell align=left]
\plotcoords{cE, thick}{\wgSwAtr}
\addlegendentry{trained}
\plotcoords{ink2, dashed}{\wgSwAbi}
\addlegendentry{Biagi's file}
\plotcoords{cH, dotted, thick}{\wgSwAHM}
\addlegendentry{trained, HM band's top}
\plotcoords{only marks, mark=*, mark size=1.1pt, ink}{\wgSwATdata}
\addlegendentry{train}
\plotcoords{only marks, mark=o, mark size=1.1pt, ink2}{\wgSwAHdata}
\addlegendentry{held out}
\end{axis}
\end{tikzpicture}\hfill
\begin{tikzpicture}
\begin{axis}[width=5.5cm, height=5.6cm, xmode=log, ymode=log,
  xlabel={$E/N$ (Td)}, ylabel={$D_T/\mu$ (V)},
  tick label style={font=\scriptsize}, label style={font=\footnotesize},
  legend style={font=\tiny, draw=none, fill=none, at={(0.03,0.97)},
  anchor=north west}, legend cell align=left]
\plotcoords{cE, thick}{\wgSwDtr}
\addlegendentry{trained}
\plotcoords{ink2, dashed}{\wgSwDbi}
\addlegendentry{Biagi's file}
\plotcoords{only marks, mark=o, mark size=1.1pt, ink2}{\wgSwDHdata}
\addlegendentry{held out}
\end{axis}
\end{tikzpicture}
\caption{Krypton's drift velocity (left), first Townsend coefficient
(middle) and characteristic energy (right) at 300\unit{K}, as the
experiments measure them: the trained cross sections (solid), Biagi's
file as published (dashed), and the measurements of
Table~\ref{tab:swdata}, filled where they train and open where they are
held out (the validation sets among them). Middle: the dotted curve is
the trained set with every excitation of the
levels at or above \Krtp's energy ionizing associatively (the
Hornbeck--Molnar band's top).}
\label{fig:swarm}
\end{figure}

\textbf{The scores.}
\begin{itemize}
\item \textbf{The train sets.} The drift velocity of Hunter et
al.~\cite{hunter88}, the most precise below 3\unit{Td}, moves from
$\chi^2$ \wg{sw.w_hunter88.bi.chi} to \wg{sw.w_hunter88.tr.chi} per
point, Nakamura's~\cite{nakamura90} from \wg{sw.w_nakamura90.bi.chi} to
\wg{sw.w_nakamura90.tr.chi} (mean \wg{sw.w_nakamura90.bi.mean} to
\wg{sw.w_nakamura90.tr.mean}\,\%). The ionization coefficients of
Jacques et al.~\cite{jacques86} and Kruithof~\cite{kruithof40}, which lay
\wg{sw.a_jacques86.bi.mean} and \wg{sw.a_kruithof40.bi.mean}\,\% from
the model, now lie \wg{sw.a_jacques86.tr.mean} and
\wg{sw.a_kruithof40.tr.mean}\,\%; Bhattacharya's~\cite{bhattacharya79},
which Biagi's file matched (\wg{sw.a_bhatta79.bi.mean}\,\%), now lies
\wg{sw.a_bhatta79.tr.mean}\,\%. Every train set ends within the
consistency threshold.
\item \textbf{The held-out drift velocities} are the sets the check
dropped. The training moves \wg{sw.wh.better} of the \wg{sw.wh.sets}
closer to the model, the two at 195\unit{K} among them, which took part
in choosing which factors train. Three move further: Pack et al.\ 1992
at 300\unit{K}~\cite{pack92} ($\chi^2$ \wg{sw.w_pack92.bi.chi} to
\wg{sw.w_pack92.tr.chi} per point), Dall'Armi et al.'s~\cite{dallarmi92}
(\wg{sw.w_dallarmi92.bi.chi} to \wg{sw.w_dallarmi92.tr.chi}) and
Kucukarpaci and Lucas's~\cite{kucukarpaci81} (\wg{sw.w_kucuk81.bi.chi}
to \wg{sw.w_kucuk81.tr.chi}, \wg{sw.w_kucuk81.tr.mean}\,\% from the
model on average). Pack et al.'s 1992 set scatters by up to
\wg{pack92.jump}\,\% against the model between points within 5\,\% of
each other in $E/N$; Bowe's~\cite{bowe60} lie
\wg{sw.w_bowe60.tr.mean}\,\% from the model, and English and
Hanna's~\cite{english53} \wg{sw.w_english53.tr.mean}\,\%.
\item \textbf{The characteristic energy} of Koizumi et
al.~\cite{koizumi86}, held out, probes the momentum transfer near the
Ramsauer minimum, which the drift velocities pin least: rms
\wg{sw.dtk.bi.rms}\,\% on Biagi's file and \wg{sw.dtk.tr.rms}\,\% on the
trained set. It chose which factors train, so it validates the training
rather than tests it.
\item \textbf{Below 50\unit{Td}}, held out by rule, the mean error moves
from \wg{sw.alow.bi.mean} to \wg{sw.alow.tr.mean}\,\%: Heylen's
from \wg{sw.a_heylen71_low.bi.mean} to \wg{sw.a_heylen71_low.tr.mean}\,\%,
Specht et al.'s direct coefficient~\cite{specht80} from
\wg{sw.a_specht80_direct.bi.mean} to
\wg{sw.a_specht80_direct.tr.mean}\,\%, Kruithof's from
\wg{sw.a_kruithof40_low.bi.mean} to \wg{sw.a_kruithof40_low.tr.mean}\,\%.
\end{itemize}

\textbf{Why the low field stays out.} At 30\unit{Td} the electrons excite
krypton \wg{low.exc_ion} times for every ionization, so anything that
ionizes an excited atom, or frees an electron at a surface with its
light, can dominate the measured current. Specht et al.\ found the
coefficient anomalously large at low field and attributed it to
photoelectron emission at the surfaces by radiating metastables, which
they say also explains why earlier measurements disagreed with transport
calculations; their direct coefficient leaves that contribution
out~\cite{specht80}. Bordage et al.\ found every published krypton set
below the low-field data and ascribe it to Penning ionization of
impurities~\cite{bordage13}. And the \wg{sw.low.sets} datasets part by up
to \x{\wg{sw.low.part}} from one another where they overlap. Training on
them would fit the cross sections to the purity and geometry of each
experiment. Pure krypton has one process there of its own, the
associative (Hornbeck--Molnar) ionization of atoms excited above
\wg{hm.eai}\unit{eV}~\cite{hornbeck51}. At the band's top it raises the
trained set's coefficient \x{\wg{hm.x14}} at 14\unit{Td}, \x{\wg{hm.x30}}
at 30\unit{Td} and \x{\wg{hm.x100}} at 100\unit{Td}, and puts the model
\wg{swhm.atr.mean}\,\% from the train sets on average, so a share that
large cannot hold where those are measured. No krypton measurement gives
the share, and it stays at 0 by default.

\textbf{What this means for the reference point.} The electrons sit at
\wg{ref.kin.en.lo} to \wg{ref.kin.en.hi}\unit{Td}. There direct
ionization makes \wm{ion.C2}\,\% of the ions, and the ionization of the
excited atoms and of \Krts{} with Penning ionization the rest, so the
low-field coefficient does not reach the answer. The
excitation rates do, and the training moves them: on Biagi's file as
published v0\_5 gives \wg{chg.kin.bi}\,\% at the reference point, and on
the trained set \wg{ref.kin.147}\,\% (\wg{chg.kin.bi.dev}\,\%); the
posterior of the factors spreads it by $\pm$\wg{unc.kin.xs.sd} points
(Section~\ref{sec:accuracy}).

\section{The electroluminescence yield}
\label{sec:el}

Mano et al.\ measured the secondary scintillation of pure krypton at
1.1\unit{bar} and room temperature: the photons an electron drifting in a
uniform field makes per centimetre, over the gas density~\cite{mano22}.
Below the field at which the electrons ionize, every photon comes from an
excitation, so the yield measures the excitation cross sections in the
swarm regime, where no electron meets another and no excited atom meets
an electron. Their fit, $Y/N = \wg{el.mano.a}\,E/N - \wg{el.mano.b}$ in
$10^{-17}$ photons per electron per atom per cm$^{-2}$ ($E/N$ in Td),
holds from its intercept at \wg{el.mano.thr}\unit{Td} to
\wg{el.mano.top}\unit{Td}. Nothing in v0\_5 or in its training saw it.

v0\_5's swarm (Section~\ref{sec:swarm}) gives the excitations per
electron and per centimetre of drift, $\sum k_\mathrm{exc}/W$, level by
level. Each excitation is then weighted by the photons the model's own
network makes of it at 1.1\unit{bar} and 296\unit{K}: a dilute seed of
each 5s level and of the 5p rung, with no drive and no electrons to speak
of, run to the end of its cascade. At the bands' defaults every seed ends
as one 147-nm photon (\wg{el.ph.1s5.default} per \lv{5} atom,
\wg{el.ph.1s4.default} per \lv{4} atom, \wg{el.ph.5p.default} per 5p
atom); with \Krts{} + Kr at its band's top, \wg{el.ph.5p.quench} per 5p
atom.

\begin{table}[H]
\centering\footnotesize
\begin{tabular}{@{}rrrrrr@{}}
\toprule
 & \multicolumn{2}{c}{$Y/N$ ($10^{-17}\unit{cm^2}$)} & \multicolumn{2}{c}{model over fit} & \\
\cmidrule(lr){2-3}\cmidrule(lr){4-5}
$E/N$ (Td) & Mano et al.'s fit & model & bands' defaults & \Krts{} + Kr at its top & 5s share (\%)\\
\midrule
\wgEl
\bottomrule
\end{tabular}
\caption{The electroluminescence yield of pure krypton: Mano et al.'s
fit~\cite{mano22} and v0\_5's swarm at 296\unit{K}, each excitation
weighted by its photons from the network's cascade at 1.1\unit{bar}; the
5s share is the excitations into the four 5s levels.}
\label{tab:el}
\end{table}

From 6 to 13.5\unit{Td} the model lies at \wg{el.lo} to \wg{el.hi} of the
fit (\wg{el.rms}\,\% rms), and with the excimer's quenching at its band's
top at \wg{el.q.lo} to \wg{el.q.hi}. Toward the fit's intercept it falls
below it (\wg{el.5} at 5\unit{Td}, \wg{el.4} at 4): there the fit's yield
is \wg{el.frac5} and \wg{el.frac4} of its value at 13.5\unit{Td}, and a
straight line's intercept is the least constrained part of a fit to a
yield that bends into its threshold. Nearly every excitation below 14\unit{Td} is a 5s
level's, so the score tests the four 5s cross sections and the trained
5s factor (\x{\wg{fac.x_1s}}) at the swarm's mean energies of
\wg{el.mean.lo} to \wg{el.mean.hi}\unit{eV}, and the
excitation-to-photon chain at a density where nothing competes with the
light.

\section{Under an electron beam}
\label{sec:beam}

Three measurements bear on the chain from excitation to the 147-nm
photon. They were made under electron beams, so for them v0\_5 runs under
a beam drive. The drive is a fixture of this note, not part of the model.
It deposits the given power as ion pairs, 5s and 5p atoms and secondary
electrons in the proportions of a fast electron's degradation in pure
krypton. That degradation is the Fowler--Spencer--Fano recursion on the
same cross sections as the model, level by level, from
\wg{beam.e0}\unit{eV}: an excitation continues the electron at $E-I$, and
an ionization shares $E-I$ equally between the two electrons. Per ion
pair it gives:
\begin{itemize}
\item $W=\wg{beam.W}\unit{eV}$;
\item \wg{beam.pp} \Krpp, from the levels above the 5s manifold;
\item \wg{beam.s} 5s atoms: \wg{beam.s.1s5} \lv{5}, \wg{beam.s.1s4}
\lv{4}, \wg{beam.s.1s3} \lv{3} and \wg{beam.s.1s2} \lv{2}, the two
resonance levels most, as a fast electron's dipole excitation favours
them (v0\_3 put the four into one \Krs);
\item \wg{beam.cascade}\unit{eV} of heat, the higher levels' excess over
the 5p rung's energy;
\item \wg{beam.eps}\unit{eV} of sub-excitation secondaries, which the
kinetic equation injects with a flat spectrum up to \wm{E.S}\unit{eV}.
\end{itemize}
That is one excitation or ion per \wg{beam.ws}\unit{eV}.

\textbf{The partition.} Saito et al.~\cite{saito03} measured scintillation
photons and electrons together under 5.49-MeV alpha particles in 1 to
10\unit{atm} of krypton. They found \wg{saito.nex} excited atoms per ion
pair, and $W_s=\wg{saito.ws}\unit{eV}$, the energy per photon when every
pair recombines. They also confirmed that each recombination gives one
photon. On the trained cross sections the drive's \wg{beam.ws}\unit{eV}
per excitation or ion is \wg{beam.ws.dev}\,\% from their $W_s$, its
excitations per pair (\wg{beam.nex}) differ from theirs by
\wg{beam.nex.dev}\,\%, and its $W$ differs by \wg{beam.W.dev}\,\% from
the \wg{saito.W}\unit{eV} their two numbers imply. On Biagi's file as
published the same degradation gives $W=\wg{beam.W.v02}\unit{eV}$
(\wg{beam.W.dev.v02}\,\%) and \wg{beam.nex.v02} excitations per pair
(\wg{beam.nex.dev.v02}\,\%): the swarm data's smaller excitation moves
both toward the measurement, which neither the training nor the choice
of its factors saw. One photon of \wg{hv147}\unit{eV} per
\wg{saito.ws}\unit{eV} is an efficiency of \wg{saito.ideal}. The model's
ceiling, one photon per excitation and per ion of the drive, is
\wg{beam.ceil}.

\textbf{Low pump.} Morozov et al.~\cite{morozov08} measured the power a
continuous 8 to 15-keV beam deposits in the gas against the VUV power
emitted into $4\pi$. For krypton they found \wg{moz.meas} $\pm$
\wg{moz.err}, flat from 500 to 1400\unit{hPa} and over a tenfold range of
beam power. At \wg{moz.hpa}\unit{hPa}, \wg{moz.w.lo} and
\wg{moz.w.hi}\unit{W\,cm^{-3}} for \wg{beam.pulse}\unit{\mu s} with the
light counted for \wg{beam.tail.lo}\unit{\mu s} after, v0\_5 gives
\wg{moz.30.kin} and \wg{moz.300.kin} (Table~\ref{tab:beam}). The error is
\wg{moz.kin.dev}\,\%, \wg{moz.kin.sig} of the measurement's standard
errors (v0\_4: \wg{moz.kin.dev.v04}\,\%). At this pump nothing competes
with the light: \wg{moz.300.kin.G2} Penning ionizations per pair, every
ion recombined (\wg{moz.300.kin.rec} per pair), \wg{moz.300.kin.redis}
excimers redissociated and \wg{moz.300.kin.ei} ionized per pair, so
nearly every \Krts{} radiates. So v0\_5 sits at its ceiling, which Saito et al.'s measurement supports. Morozov
et al.\ compare their value with calculated maxima of 0.46 and 0.52 for
krypton. The \wg{moz.short}\,\% between their value and Saito et al.'s
one photon per excitation and per ion (\wg{saito.ideal}) is not explained
in v0\_5.

\textbf{High pump.} Eckstrom et al.~\cite{eckstrom88} report the
fluorescence efficiency of \Krts{} near their theoretical limit of about
46\,\%. They found it under a low-current, long-pulse beam at modest
pressures, for pump rates up to $3\times10^5\unit{W\,cm^{-3}}$. The
abstract does not give the pressure of each point. At that pump rate
(1\unit{\mu s}, the light counted for \wg{beam.tail.hi}\unit{\mu s})
v0\_5 gives:
\begin{itemize}
\item \wg{eck.0.5.kin} at 0.5\unit{atm}, \wg{eck.1.kin} at 1\unit{atm}
and \wg{eck.2.kin} at 2\unit{atm}, that is \wg{eck.0.5.kin.loss},
\wg{eck.1.kin.loss} and \wg{eck.2.kin.loss}\,\% below its ceiling of
\wg{beam.ceil} (v0\_4: \wg{eck.0.5.kin.v04}, \wg{eck.1.kin.v04} and
\wg{eck.2.kin.v04});
\item at 1\unit{atm}, Penning ionization of the 5s levels on one another
costs \wg{eck.1.kin.G2} photons per pair, because the excimer forms too
slowly to outrun it entirely;
\item at 1\unit{atm} with \Krs{} + \Krs{} at its band's low end,
\wg{eck_pen.1.kin} (\wg{eck_pen.1.kin.loss}\,\% below the ceiling), and
with the \lv{5} formation at its band's top, \wg{eck_form.1.kin}
(\wg{eck_form.1.kin.loss}\,\%).
\end{itemize}
The rows v0\_5 adds act under the beam only where its plasma is dense. At
low pump the efficiency is v0\_4's, \wg{moz.kin.v04}. At
$3\times10^5\unit{W\,cm^{-3}}$ and 1\unit{atm} the electrons ionize
\wg{eck.1.kin.ei} excimers per ion pair and the stepwise ionizations make
\wg{eck.1.kin.step} more ions, and the light changes by
\wg{eck.1.vs.v04}\,\% against v0\_4; at 0.5\unit{atm}, where the gas
passes \wg{eck.0.5.kin.tg}\unit{K}, by \wg{eck.0.5.vs.v04}\,\%, and at
2\unit{atm} by \wg{eck.2.vs.v04}\,\%.

\begin{table}[H]
\centering\footnotesize
\resizebox{\linewidth}{!}{%
\begin{tabular}{@{}llrrrrrrr@{}}
\toprule
run & measured & kinetic & two-temperature & v0\_4, kinetic & Penning per pair & photons per pair & $T_e$ (eV) & $T_g$ end (K)\\
\midrule
\wg{moz.w.lo}\unit{W\,cm^{-3}}, \wg{moz.hpa}\unit{hPa} & \wg{moz.meas} $\pm$ \wg{moz.err} & \wg{moz.30.kin} & \wg{moz.30.2t} & & \wg{moz.30.kin.G2} & \wg{moz.30.kin.ph} & \wg{moz.30.kin.te} & \wg{moz.30.kin.tg}\\
\wg{moz.w.hi}\unit{W\,cm^{-3}}, \wg{moz.hpa}\unit{hPa} & \wg{moz.meas} $\pm$ \wg{moz.err} & \wg{moz.300.kin} & \wg{moz.300.2t} & \wg{moz.kin.v04} & \wg{moz.300.kin.G2} & \wg{moz.300.kin.ph} & \wg{moz.300.kin.te} & \wg{moz.300.kin.tg}\\
$3\times10^5\unit{W\,cm^{-3}}$, 0.5\unit{atm} & near \wg{eck.meas} & \wg{eck.0.5.kin} & \wg{eck.0.5.2t} & \wg{eck.0.5.kin.v04} & \wg{eck.0.5.kin.G2} & \wg{eck.0.5.kin.ph} & \wg{eck.0.5.kin.te} & \wg{eck.0.5.kin.tg}\\
$3\times10^5\unit{W\,cm^{-3}}$, 1\unit{atm} & near \wg{eck.meas} & \wg{eck.1.kin} & \wg{eck.1.2t} & \wg{eck.1.kin.v04} & \wg{eck.1.kin.G2} & \wg{eck.1.kin.ph} & \wg{eck.1.kin.te} & \wg{eck.1.kin.tg}\\
$3\times10^5\unit{W\,cm^{-3}}$, 2\unit{atm} & near \wg{eck.meas} & \wg{eck.2.kin} & \wg{eck.2.2t} & \wg{eck.2.kin.v04} & \wg{eck.2.kin.G2} & \wg{eck.2.kin.ph} & \wg{eck.2.kin.te} & \wg{eck.2.kin.tg}\\
\quad 1\unit{atm}, \Krs{} + \Krs{} $\wm{bd.kr_pen.lo}$ & & \wg{eck_pen.1.kin} & \wg{eck_pen.1.2t} & & \wg{eck_pen.1.kin.G2} & \wg{eck_pen.1.kin.ph} & \wg{eck_pen.1.kin.te} & \wg{eck_pen.1.kin.tg}\\
\quad 1\unit{atm}, \lv{5} + 2Kr $\wm{bd.kr2_form.hi}$ & & \wg{eck_form.1.kin} & \wg{eck_form.1.2t} & & \wg{eck_form.1.kin.G2} & \wg{eck_form.1.kin.ph} & \wg{eck_form.1.kin.te} & \wg{eck_form.1.kin.tg}\\
\bottomrule
\end{tabular}}
\caption{v0\_5 under the beam drive: the VUV efficiency (147-nm energy
over the deposited energy) on both closures and v0\_4's on the kinetic
EEDF, and on the kinetic EEDF the Penning ionizations and 147-nm photons
per ion pair, the pulse's mean $T_e$ and the final gas temperature. The
ceiling, one photon per excitation and per ion, is \wg{beam.php} photons
per pair, an efficiency of \wg{beam.ceil}; Saito et al.'s ideal is
\wg{saito.ideal}.}
\label{tab:beam}
\end{table}

\section{The energy ladder}

Figure~\ref{fig:ladder} draws v0\_5's levels at their formation energies
and the rows that move energy between them.

\begin{figure}[H]
\centering
\begin{tikzpicture}[x=1cm, y=0.6cm, font=\scriptsize,
  arr/.style={-{Stealth[length=4pt]}},
  heat/.style={cH, decorate, decoration={snake, amplitude=1pt,
    segment length=4pt, post length=4pt}, -{Stealth[length=4pt]}}]
\newcommand\lev[4]{\draw[line width=1.1pt, #4] (#1,#3) -- (#2,#3);}
% atoms: the four 5s levels, the 5p rung, the ion
\lev{0}{4.6}{0}{ink}
\lev{0}{1.15}{\wg{lad.1s5}}{cX}
\lev{1.15}{2.3}{\wg{lad.1s4}}{cX}
\lev{2.3}{3.45}{\wg{lad.1s3}}{cX}
\lev{3.45}{4.6}{\wg{lad.1s2}}{cX}
\lev{0}{4.6}{\wg{lad.P}}{cP}
\draw[cP, dashed] (0,\wg{lad.rung.top}) -- (4.6,\wg{lad.rung.top});
\lev{0}{4.6}{\wg{lad.I}}{cI}
\node[left] at (0,0) {Kr, 0};
\node[left, align=right] at (0,\wg{lad.1s5}+0.2) {\Krs{} (5s): \lv{5}
  \wg{lad.1s5},\\ \lv{4} \wg{lad.1s4}, \lv{3} \wg{lad.1s3}, \lv{2}
  \wg{lad.1s2}};
\node[left] at (0,\wg{lad.P}) {\Krpp{} (5p) \wg{lad.P}};
\node[left, cP] at (0,\wg{lad.rung.top}+0.25) {its levels to \wg{lad.rung.top}};
\node[left] at (0,\wg{lad.I}) {\Krp{} \wg{lad.I}};
% molecules
\lev{7.4}{9.6}{0}{ink}
\lev{7.4}{9.6}{\wg{lad.X2}}{cV}
\lev{7.4}{9.6}{\wg{lad.I2}}{cI}
\draw[ink2, dashed] (7.4,\wg{lad.low}) -- (9.6,\wg{lad.low});
\node[right] at (9.6,0) {2Kr, 0};
\node[right, align=left] at (9.6,\wg{lad.X2}) {\Krts{} \wg{lad.X2}\\
  (0$_u^+$ \wg{lad.F2})};
\node[right] at (9.6,\wg{lad.I2}) {\Krtp{} \wg{lad.I2}};
\node[right, ink2] at (9.6,\wg{lad.low}) {repulsive 2Kr, \wg{lad.low}};
% electron impact from the ground state
\draw[cE, arr] (0.4,0) -- (0.4,\wg{lad.1s5}) node[pos=0.45, left] {G1};
\draw[cE, arr] (1.0,0) -- (1.0,\wg{lad.P}) node[pos=0.45, right] {C1};
\draw[cE, arr] (1.6,0) -- (1.6,\wg{lad.I}) node[pos=0.45, right] {C2};
% Penning ionization of the 5s levels on one another
\draw[ink2, densely dashed, arr] (2.3,\wg{lad.1s3}) -- (2.3,\wg{lad.I})
  node[pos=0.75, right] {G2};
% the 5s <-> 5p cycle, its relaxation and its light
\draw[cE, arr] (3.0,\wg{lad.1s3}) -- (3.0,\wg{lad.P});
\draw[cE, arr] (3.25,\wg{lad.P}) -- (3.25,\wg{lad.1s3});
\draw[heat] (3.75,\wg{lad.P}) -- (3.75,\wg{lad.1s2});
\draw[cP, arr] (4.25,\wg{lad.P}) -- (4.25,\wg{lad.1s2});
\node[ink, anchor=north west, align=left, inner sep=1pt] at (2.75,9.0)
  {\textcolor{cE}{K76, K76r,}\\ \textcolor{cH}{KPS},
  \textcolor{cP}{PRAD}\\ \textcolor{ink2}{each 5s level's}};
% v0_5: the electrons ionize the excited atoms and return the 5s levels
\draw[cE, arr] (4.45,\wg{lad.1s2}) -- (4.45,\wg{lad.I});
\node[cE, anchor=east] at (4.4,\wg{lad.P}+0.5) {KSTEP};
\draw[cE, arr, densely dotted] (0.7,\wg{lad.1s5}) -- (0.7,0)
  node[midway, sloped, above=-1pt, font=\tiny, cE] {SUPER};
% to the excimer, its reverse and its light
\draw[cX, arr] (4.6,\wg{lad.1s5}) -- (7.4,\wg{lad.X2});
\draw[cX, arr, densely dotted] (7.4,\wg{lad.X2}+0.45) -- (4.6,\wg{lad.1s5}+0.45);
\node[cX, anchor=north west, align=left] at (4.7,8.65) {AK8, \lv{5}:
  \wg{lad.ak8} to heat;\\ AK8r back, dotted};
\draw[cV, arr] (8.1,\wg{lad.X2}) -- (8.1,\wg{lad.low})
  node[pos=0.5, right, cV] {KR2\_RAD, $h\nu$ \wg{lad.hv}};
\draw[heat] (8.1,\wg{lad.low}) -- (8.1,0);
% v0_5: the electrons ionize the excimer
\draw[cE, arr] (9.35,\wg{lad.X2}) -- (9.35,\wg{lad.I2})
  node[pos=0.4, right, cE] {EI};
% the ions
\draw[cI, arr] (4.6,\wg{lad.I}) -- (7.4,\wg{lad.I2});
\node[cI, anchor=south west] at (7.45,13.25) {K15 (K15r back), \wg{lad.k15} to heat};
\draw[cI, arr] (7.4,\wg{lad.I2}) -- (4.6,\wg{lad.P});
\node[cI, anchor=north west, align=left] at (7.45,12.4) {R11p,\\ \wg{lad.r11} to heat};
\end{tikzpicture}
\caption{v0\_5's energy ladder in eV above ground-state krypton at rest.
A bar is a species at its formation energy; the four 5s levels are
drawn side by side. Blue arrows are electron impact, which takes its
energy from the electron gas; K76r, the 5p \ra 5s arrow beside K76,
returns it; v0\_5's KSTEP (KSTEP\_S and KSTEP\_P) ionize the 5s levels and the
rung (\wg{lad.ion.1s2} to \wg{lad.ion.1s5}\unit{eV} and
\wg{lad.stepP}\unit{eV} from 2p$_{10}$), EI the excimer
(\wg{lad.ei}\unit{eV}), and SUPER (dotted) returns a 5s level to the
ground state, its energy to the electron. The dashed bar is the rung's
top level; its levels are populated at the gas temperature, and a pump
to one above 2p$_{10}$ heats the gas with the excess as they relax. The
dashed G2 is the Penning ionization of the 5s levels on
one another, which gives the electron \wg{lad.pen}\unit{eV} for two
\lv{5} atoms. A red squiggle is heat: KPS relaxes the 5p rung on
krypton, and PRAD radiates its gap to each level, \wg{lad.gap.1s2} to
\wg{lad.gap.1s5}\unit{eV}, in the near infrared. The excimer forms out
of each 5s level, its well and the level's excess to heat, redissociates
by detailed balance (AK8r), and radiates onto the repulsive ground state,
which falls apart and heats the gas. The 5s levels exchange population on
krypton (MIX, not drawn). C1's levels, \wg{lad.c1lo} to
\wg{lad.c1hi}\unit{eV}, cascade to the 5p rung, each heating the gas with
its excess. HM, the Hornbeck--Molnar ionization of the levels above
\Krtp's energy, and the resonance lines of \lv{4} and \lv{2} are bands at
0 and not drawn.}
\label{fig:ladder}
\end{figure}

\section{Where v0\_5 is not yet reliable}
\label{sec:limits}

\begin{itemize}
\item \textbf{The excimer's formation.} Its rates still decide much of
the answer. The \lv{5} rate's measurements spread by a factor of
\wg{form.ratio}; the \lv{4} rate is one measurement~\cite{marchal18},
which Audouard et al.'s formation-time ratio puts three times
lower~\cite{audouard88} and Janssens et al.'s quasiclassical calculation
1.5 times higher~\cite{janssens87}; the \lv{3} rate is a proxy, Kholin et
al.'s measurement having argon for its third body~\cite{kholin15}. The
temperature law is measured over 200 to 400\unit{K}~\cite{wieme81}, and
the reference point ends at \wg{ref.kin.tg}\unit{K}.
\item \textbf{The calculated rows.} The stepwise ionizations of the 5s
levels and of the 5p rung and the ionization of \Krts{} by electrons run
on Vriens and Smeets's calculation~\cite{vriens80}, and no krypton
measurement tests it. At the reference point they make
\wg{ref.kin.ion.step_s}, \wg{ref.kin.ion.step_p} and
\wg{ref.kin.ion.ei}\,\% of the ions, and their bands, 0 to twice the
calculation, take the 147-nm share from \wg{unc.end.kstep_s.low} to
\wg{unc.end.kstep_s.high}\,\% (5s), \wg{unc.end.kstep_p.low} to
\wg{unc.end.kstep_p.high}\,\% (5p) and \wg{unc.end.ei_kr2.low} to
\wg{unc.end.ei_kr2.high}\,\% (\Krts). Kannari et al.'s fits of the same
rows~\cite{kannari85} lie \wg{kann.lo} to \wg{kann.hi} times the
calculation for $T_e$ of 1 to 3\unit{eV}, and the calculation's peak for
\Krts, $\wg{vs.kr2.peak}\unit{cm^2}$, is the order of Flannery et al.'s
binary-encounter cross sections~\cite{flannery81}.
\item \textbf{The electrons' other collisions with \Krts}, its quenching
to the ground state and the mixing of its two states, are not measured;
the one value found is Kannari et al.'s estimate of the quenching, the
top of both bands~\cite{kannari85}. Both are 0 by default; at the top the
quenching takes the share to \wg{unc.end.eq_kr2.high}\,\% and the mixing
to \wg{unc.end.em_kr2.high}\,\%.
\item \textbf{The superelastic returns} of the 5s levels are the
detailed-balance reverses of Biagi's excitations. For thermal electrons
the model's rates are \wg{blag.1s5.lo} to \wg{blag.1s5.hi} times Blagoev
et al.'s~\cite{blagoev84} for \lv{5}, \wg{blag.1s4.lo} to
\wg{blag.1s4.hi} times for \lv{4} and \wg{blag.1s3.lo} to
\wg{blag.1s3.hi} times for \lv{3} (Section~\ref{sec:excited}); such
electrons read the excitations within a few tenths of an eV of their
thresholds. At the reference point the returns move the share by
\wg{undo.no_super.kin.back}\,\%.
\item \textbf{The 5s levels' mixing by electrons} carries \lv{5}
$\leftrightarrow$ \lv{4} alone, on Blagoev and Pilosov's afterglow
rate~\cite{blagoev81}, measured at an electron temperature their
abstract does not give; the band's top takes it as measured at
0.07\unit{eV}. Qiu et
al.\ measured the \lv{5} \ra \lv{4} and \lv{3} \ra \lv{2} transfers for
$T_e$ of 0.07 to 1\unit{eV}~\cite{qiu18}, values not in a source in
hand; the other pairs are not carried.
\item \textbf{The 5p rung's levels} are populated at the gas
temperature, as if collisions with krypton mixed them faster than they
radiate, are pumped or relax. That sets which level a pump's reverse, a
stepwise ionization and a line leave from. Whitehead et al.\ measured the
deactivation of the two-photon-excited 5p levels in
krypton~\cite{whitehead95}, values not in a source in hand.
\item \textbf{The field-driven regime is unscored.} No pure-krypton
measurement in hand tests the 5s $\leftrightarrow$ 5p cycle under a field.
Its cross sections are measured to $\pm$\wg{jung.unc.1s5} to
\wg{jung.unc.1s3}\,\% out of the two metastables; those of the two $J=1$
levels are not measured and take their metastable's (band
\texttt{k\_5s5p\_j1}). Zeng et al.'s relativistic R-matrix calculation
gives them~\cite{zeng05}, its values not in a source in hand.
\item \textbf{The redissociation} rests on the excimer's well, whose
measurements give 0.62 to 0.709\unit{eV} (its band), and on a rotational
constant from an estimated bond length, 2.6\,\AA, the ion's, which enters
as $1/r_e^2$; the ion's vibrational quantum is taken as the excimer's. At
the reference point \wg{ref.kin.redis}\,\% of the excimers formed
redissociate; above \wg{keq.half.K}\unit{K} at the fill's density most
would.
\item \textbf{Hot gas.} v0\_5 carries the rung's reverse, which reaches
\wg{kpsr.max}\,\% of its relaxation at the scan's \wg{kpsr.max.atm}\unit{atm}
and \wg{kpsr.max.kw}\unit{kW\,cm^{-3}}. The scan's 0.1-atm column ends at
\wg{scan.tg01.lo} to \wg{scan.tg01.hi}\unit{K}, outside the measured
ranges of the formation and the recombination.
\item \textbf{The excimer's two states.} The fast 0$_u^+$ state takes no
share by default: its measured share is above $6\times10^{19}\unit{cm^{-3}}$
and appears from $1.5\times10^{19}$~\cite{audouard88}, so at 1\unit{atm}
it lies between 0 and the band's top. Quenching by krypton is 0 by
default, against one measurement of a pressure-dependent decay at
984\unit{nm}~\cite{firestone88} that may be a transfer between radiating
states rather than a loss. No rate mixes the two states on krypton, and
the electrons' mixing is a band at 0.
\item \textbf{The two-temperature closure} now reads its tail at the
run's excited fraction. With that axis it lies \wg{tail.ref.f.dev}\,\%
from the kinetic EEDF at the reference point (on v0\_4's table, without
the axis, \wg{tail.ref.f0.dev}\,\%) and within \wg{tail.worst.f}\,\% at the
scan's corners. Its tail now follows the kinetic EEDF's along the reference run
(\wg{tail.f.lo} to \wg{tail.f.hi} of it), so what is left is its
Maxwellian below the break, which v0\_4's table had more than offset.
The default closure does not use it.
\item \textbf{The trained cross sections} rest on errors the swarm data
do not print (class values here) and on the consistency check's threshold
(at 4 rather than 3, \wg{cons.at4.n} more sets train and the momentum
transfer at 6\unit{eV} moves to \x{\wg{cons.at4.m6}}). The data barely
pin the 5p excitation. Which factors train was chosen on held-out data
(Table~\ref{tab:var}), so those sets validate the training rather than
test it. Their posterior is propagated into the answer
($\pm$\wg{unc.kin.xs.sd} points).
\item \textbf{The rows with no measurement or calculation.} \Krts{} +
\Krts{}, the electrons' quenching and mixing of \Krts{} and the
Hornbeck--Molnar share are 0 by default (Table~\ref{tab:unc}). At its
band's top \Krts{} + \Krts{} moves the 147-nm share by
\wg{unc.high.kr2_pen}\,\%. The $J=1$ levels' 5s \ra 5p cross sections
are proxies. With the excimer's ionization, the two rows on \Krts{} at 0 make the
uncertainty's three largest terms (Section~\ref{sec:accuracy}).
\item \textbf{\Krs{} + \Krs.} Its $\wm{k.G2}\cm{3}$ is a 1974
measurement~\cite{leichner74} as Kannari et al.\ print it~\cite{kannari85},
taken for every pair of levels.
\item \textbf{The e-beam efficiency at low pump} is off Morozov et
al.'s by \wg{moz.kin.dev}\,\%, and v0\_5 holds no process that would
explain the difference.
\item \textbf{The ionization coefficient below 50\unit{Td}} lies
\wg{sw.alow.tr.mean}\,\% from the measurements on average, which part by
up to \x{\wg{sw.low.part}} from one another; every published
cross-section set falls below them~\cite{bordage13}. It is held out of the
training by rule, and it does not reach the reference point, where the
ionization of the excited atoms and of \Krts{} makes the ions.
\item \textbf{The Hornbeck--Molnar share} of the levels above
\wg{hm.eai}\unit{eV} is not measured for krypton and is 0 by default. At
its top it puts the model \wg{swhm.atr.mean}\,\% from the train sets and
moves the reference point by \wg{unc.hm.cut}\,\%.
\item \textbf{Processes v0\_5 does not carry}, each needing a measured
rate or a calculation before it can enter:
\begin{itemize}
\item Kr$_3^+$ and its recombination;
\item Penning ionization between \Krs{} and \Krts, and of \Krpp;
\item the electrons' excitation of the 5s and 5p levels above the rung,
their transfers between the 5s levels other than \lv{5} $\leftrightarrow$
\lv{4}, and their return of the rung to the ground state.
\end{itemize}
Three-body recombination \Krp{} + \el{} + \el~\cite{mansbach69} is not
carried either: at 1\unit{atm} \Krp{} converts to \Krtp{} in
\wg{k15.ns}\unit{ns}.
\end{itemize}

\textbf{How the numbers were made.} Every number in this note is computed
from the model when the note is built. None is typed by hand except the
measurements, which carry their sources, and the training's record, which
the model carries beside its cross sections and the training wrote. The
runs are cached under a hash of the model's code, so a change to the
model reruns them all (code \texttt{\wg{md5}}). The kinetic runs take
about \wg{ref.kin.min} minutes each; the two-temperature runs take
seconds.

\clearpage
\appendix
\section{Reaction rates}
\label{app:rates}

The rows of v0\_5 with the values they run with and their sources, a
\texttt{<level>} or \texttt{<1s>} row once per 5s level and a
\texttt{<2p>} row once per level of the 5p rung. Electron-impact rows use
their cross sections over the kinetic EEDF (or the two-temperature
closure). A band gives its default and its two ends.

\begin{table}[H]
\centering\scriptsize
\renewcommand{\arraystretch}{1.1}
\begin{tabularx}{\linewidth}{@{}>{\raggedright\arraybackslash}p{2.45cm}>{\raggedright\arraybackslash}p{3.1cm}>{\raggedright\arraybackslash}p{4.75cm}>{\raggedright\arraybackslash}X@{}}
\toprule
row & reaction & value & source\\
\midrule
G1\_<level> & \el{} + Kr \ra \el{} + \Krs(level) & each 5s level at its threshold, \wm{E.1s5} to \wg{lad.s2}\unit{eV}, its own species & Biagi, Magboltz 8.97~\cite{biagi99,pancheshnyi12,carbone21}, \x{\wg{fac.x_1s}} (trained, Table~\ref{tab:fac})\\
C1 & \el{} + Kr \ra \el{} + \Krpp & the 5p, 4d, 6s and higher levels, \wg{lad.c1lo} to \wg{lad.c1hi}\unit{eV}, \x{\wg{fac.x_2p}} (5p) and \x{\wg{fac.x_hi}} (higher) trained; the excess over \wm{E.PP} to heat & Biagi, Magboltz 8.97\\
HM & \el{} + Kr (+ Kr) \ra \Krtp{} + 2\el & the share of the excitations of the \wg{hm.levels} levels at or above \wg{hm.eai}\unit{eV}; band 0 (default) to 1 & no measurement for krypton; the process, Hornbeck and Molnar~\cite{hornbeck51}\\
C2 & \el{} + Kr \ra \Krp{} + 2\el & \wm{E.ion}\unit{eV} & Biagi, Magboltz 8.97\\
G6 & \el{} + Kr, elastic & the momentum transfer cross section, times the trained factor & Biagi, Magboltz 8.97\\
K76\_<1s>\_<2p>, K76r & \el{} + \Krs(1s) $\rightleftarrows$ \el{} + \Krpp(2p) & the measured cross section of each of the \wm{nPairs} pairs out of \lv{5} and \lv{3}, the $J=1$ levels on their same-core metastable's, each at its own threshold on Biagi's level energies, the level's excess over 2p$_{10}$ to heat; band \x{} \wm{bd.k_5s5p.lo} to \wm{bd.k_5s5p.hi}, and \x{} 0 to 1 on the $J=1$ levels; the reverse by detailed balance out of the rung's level & Jung et al.~\cite{jung06,jung05}; Kannari's fit as a comparison~\cite{kannari85,jacob76}\\
KSTEP\_S\_<level>, KSTEP\_P\_<2p> & \el{} + \Krs(level), \Krpp(2p) \ra \Krp{} + 2\el & Vriens and Smeets's cross section at the level's ionization energy, \wm{E.stepS}\unit{eV} out of \lv{5} and \wm{E.stepP}\unit{eV} out of 2p$_{10}$, the rung's levels at $T_g$; bands \x{} 1, \wm{bd.kstep_s.lo} to \wm{bd.kstep_s.hi} & calculation~\cite{vriens80}, no measurement; Kannari et al.'s fits of Vriens's calculation~\cite{kannari85,vriens65} as a comparison\\
EI\_KR2, EI\_KR2F & \el{} + \Krts, \Krts(0$_u^+$) \ra \Krtp{} + 2\el & Vriens and Smeets's cross section at the excimer's ionization energy, \wm{E.ei}\unit{eV}; band \x{} 1, \wm{bd.ei_kr2.lo} to \wm{bd.ei_kr2.hi} & calculation~\cite{vriens80}; Flannery et al.'s order~\cite{flannery81}; Kannari et al.'s fit~\cite{kannari85}\\
EQ\_KR2 & \el{} + \Krts{} \ra \el{} + 2Kr & 0; band to $\wm{bd.eq_kr2.hi}\cm{3}$; the electron gains the 147-nm quantum & no measurement; Kannari et al.'s estimate is the top~\cite{kannari85}\\
EM\_KR2, EM\_KR2r & \el{} + \Krts{} $\rightleftarrows$ \el{} + \Krts(0$_u^+$) & 0; band to $\wm{bd.em_kr2.hi}\cm{3}$ on $e^{-E_\mathrm{g}/T_e}$, $E_\mathrm{g}$ the states' gap, \wg{lad.mix2}\unit{eV}; the reverse by detailed balance & no measurement\\
SUPER\_<level> & \el{} + \Krs(level) \ra \el{} + Kr & G1's cross section reversed by detailed balance; the electron gains the level's energy & Biagi, Magboltz 8.97, trained; scored on Blagoev et al.'s thermal rates~\cite{blagoev84}\\
EMIX\_1s5, EMIXr\_1s4 & \el{} + \Krs(\lv{5}) $\rightleftarrows$ \el{} + \Krs(\lv{4}) & $\wm{bd.emix_5s.def}\,e^{-E_\mathrm{g}/T_e}\cm{3}$, $E_\mathrm{g}$ the levels' gap, \wg{lad.mix15}\unit{eV}; band to $\wm{bd.emix_5s.hi}$, the measurement taken at 0.07\unit{eV}; the reverse by detailed balance & Blagoev and Pilosov~\cite{blagoev81}; Qiu et al.'s range of $T_e$~\cite{qiu18}\\
R11p & \Krtp{} + \el{} \ra \Krpp{} + Kr & $\wm{k.R11p}\,(300\,\mathrm{K}/T_e)^{\wm{pow.R11p}}\cm{3}$, $\pm$10\,\%; a gas-temperature factor, band 0 to 1 & Shiu and Biondi~\cite{shiu77}; Mikus within 8\,\%~\cite{mikus78}; Cunningham and Hobson~\cite{cunningham72}\\
\bottomrule
\end{tabularx}
\caption{v0\_5's electron collisions. Rates in cm$^3$\,s$^{-1}$, $T_e$ in eV.}
\label{tab:rates_e}
\end{table}

\begin{table}[H]
\centering\scriptsize
\renewcommand{\arraystretch}{1.1}
\begin{tabularx}{\linewidth}{@{}>{\raggedright\arraybackslash}p{2.45cm}>{\raggedright\arraybackslash}p{3.1cm}>{\raggedright\arraybackslash}p{4.75cm}>{\raggedright\arraybackslash}X@{}}
\toprule
row & reaction & value & source\\
\midrule
KPS\_<level>, KPSr & \Krpp{} + Kr $\rightleftarrows$ \Krs(level) + Kr & $\wm{k.KPS}\cm{3}$ in all, over the levels by $2J+1$ (band \texttt{kps\_5s}: 1, to all into \lv{5}); band $\wm{bd.kps_kr.lo}$ to the default; the reverse by detailed balance on the rung's levels at $T_g$ & Gazeli et al.~\cite{gazeli21}, one level, products unresolved; the low end, the range Chang and Setser measured for argon's 4p levels~\cite{chang78}\\
PRAD\_<2p>\_<1s> & \Krpp(2p) \ra \Krs(1s) + $h\nu$, near IR & NIST's Einstein coefficient of each of the \wm{nLines} lines, out of the rung's levels at $T_g$: at 300\unit{K} the rung's lifetime is \wm{tauPRAD}\unit{ns}, nearly all 2p$_{10}$'s (v0\_4: \wg{tau5p}\unit{ns}, the $2J+1$ mean) & NIST ASD~\cite{nistasd}; the measured lifetimes~\cite{dzierzega00} and a TALIF lifetime of one level, 34.1\unit{ns}~\cite{gazeli21rev}, as cross-checks\\
G2\_<a>\_<b> & \Krs(a) + \Krs(b) \ra \Krp{} + Kr + \el & $\wm{k.G2}\cm{3}$ for each pair (twice for two levels); band $\wm{bd.kr_pen.lo}$ to the default & Leichner and Ericson~\cite{leichner74} as Kannari et al.\ print it~\cite{kannari85}\\
AK8\_1s5 & \Krs(\lv{5}) + 2Kr \ra \Krts{} + Kr & $\wm{k.AK8}\cm{6}\times(T_g/300)^{\wm{tpow.AK8}}$; band $\wm{bd.kr2_form.lo}$ to $\wm{bd.kr2_form.hi}$ & Marchal et al.~\cite{marchal18}; Ueda et al.~\cite{ueda84}, Tracy and Oskam~\cite{tracy76}, Turner~\cite{turner67} and Smith and Turner~\cite{smith63} span the band; the law, Wieme and Lenaerts, 200 to 400\unit{K}~\cite{wieme81}\\
AK8\_1s4, AK8f & \Krs(\lv{4}) + 2Kr \ra \Krts{} + Kr & $\wm{k.AK8.1s4}\cm{6}$, the same law; band to a third of the \lv{5} rate; the share into the fast 0$_u^+$ state, band \wm{bd.kr2_fast.lo} to \wm{bd.kr2_fast.hi} & Marchal et al.~\cite{marchal18}; Audouard et al.~\cite{audouard88}\\
AK8\_1s3 & \Krs(\lv{3}) + 2Kr \ra \Krts{} + Kr & the \lv{5} rate as a proxy; band 0 to the \lv{4} rate & Kholin et al.'s measurement has argon for its third body~\cite{kholin15}\\
AK8r\_<level> & \Krts{} + Kr \ra \Krs(level) + 2Kr & detailed balance on the excimer's well, $D_0$ \wm{bd.kr2_d0.def}\unit{eV}; band to \wm{bd.kr2_d0.lo}\unit{eV} & LaRocque et al.~\cite{larocque86}, Killeen and Eden~\cite{killeen86}\\
MIX & \Krs(\lv{4}) + Kr \ra \Krs(\lv{5}) + Kr, and \Krs(\lv{2}) + 2Kr \ra \Krs(\lv{4}) + 2Kr & $1.3\times10^{-13}\cm{3}$ at 295\unit{K} on $(T_g/295)^{-1.29}$ through Turner's 196-K point, and $0.2\times10^{-32}\cm{6}$; their reverses by detailed balance & Turner~\cite{turner67}; Audouard et al.~\cite{audouard91}\\
RES & \Krs(\lv{4}, \lv{2}) \ra Kr + $h\nu$, 123.6 and 116.5\unit{nm} & the escaping share; band 0 to $1.7\times10^5\unit{s^{-1}}$ & the pressure-independent term of Janssens et al.'s loss rate~\cite{janssens87}\\
KR2\_RAD & \Krts{} \ra 2Kr + $h\nu$, 147\unit{nm} & $\tau=\wm{tauRAD}\unit{ns}$, $h\nu=\wm{hv.RAD}\unit{eV}$ & Morikawa et al., 270 $\pm$ 5\unit{ns}~\cite{morikawa89}; 250 $\pm$ 20~\cite{salamero79}, 280 $\pm$ 30~\cite{quigley78}\\
KR2F\_RAD & \Krts(0$_u^+$) \ra 2Kr + $h\nu$ & $\tau=\wm{tauRADF}\unit{ns}$ & Morikawa et al.~\cite{morikawa89}\\
QKR2 & \Krts{} + Kr \ra 3Kr & 0; band to $\wm{bd.kr2_quench.hi}\cm{3}$ & Oka et al.~\cite{oka74}; Firestone and Selander~\cite{firestone88}\\
PEN\_KR2 & \Krts{} + \Krts{} \ra \Krtp{} + 2Kr + \el & 0; band to $\wm{bd.kr2_pen.hi}$ & no measurement; Kannari et al.'s value is the top~\cite{kannari85}\\
K15, K15r & \Krp{} + 2Kr $\rightleftarrows$ \Krtp{} + Kr & $\wm{k.K15}\cm{6}\times(T_g/300)^{\wm{tpow.K15}}$; the reverse by detailed balance on the ion's \wg{lad.k15}\unit{eV} & the mean of three measurements~\cite{bhattacharya70,grossl81,smith72}; the law, P\'erez-R\'ios and Greene~\cite{perezrios18}; Signorell et al.~\cite{signorell01}\\
\bottomrule
\end{tabularx}
\caption{v0\_5's heavy-particle rows and radiation. Rates in cm$^3$\,s$^{-1}$ or cm$^6$\,s$^{-1}$. \Krtp{} lies \wg{lad.k15}\unit{eV} below \Krp{} + Kr~\cite{signorell01}.}
\label{tab:rates_h}
\end{table}

\begin{thebibliography}{99}
\footnotesize
\setlength{\itemsep}{1pt}
\bibitem{biagi99} S.~F.~Biagi, Nucl.\ Instrum.\ Methods A 421, 234 (1999), doi:10.1016/S0168-9002(98)01233-9.
\bibitem{bordage13} M.~C.~Bordage, S.~F.~Biagi, L.~L.~Alves, K.~Bartschat, S.~Chowdhury, L.~C.~Pitchford, G.~J.~M.~Hagelaar, W.~L.~Morgan, V.~Puech and O.~Zatsarinny, J.\ Phys.\ D 46, 334003 (2013), doi:10.1088/0022-3727/46/33/334003.
\bibitem{dias23} T.~C.~Dias, A.~Tejero-del-Caz, L.~L.~Alves and V.~Guerra, Comput.\ Phys.\ Commun.\ 282, 108554 (2023), doi:10.1016/j.cpc.2022.108554.
\bibitem{hornbeck51} J.~A.~Hornbeck and J.~P.~Molnar, Phys.\ Rev.\ 84, 621 (1951), doi:10.1103/PhysRev.84.621.
\bibitem{jacques86} L.~Jacques, W.~Bruynooghe, R.~Boucique and W.~Wieme, J.\ Phys.\ D 19, 1731 (1986), doi:10.1088/0022-3727/19/9/017.
\bibitem{pancheshnyi12} S.~Pancheshnyi et al., Chem.\ Phys.\ 398, 148 (2012), doi:10.1016/j.chemphys.2011.04.020.
\bibitem{carbone21} E.~Carbone et al., Atoms 9, 16 (2021), doi:10.3390/atoms9010016.
\bibitem{hagelaar05} G.~J.~M.~Hagelaar and L.~C.~Pitchford, Plasma Sources Sci.\ Technol.\ 14, 722 (2005), doi:10.1088/0963-0252/14/4/011.
\bibitem{jung05} R.~O.~Jung, T.~E.~Stone, J.~B.~Boffard, L.~W.~Anderson and C.~C.~Lin, Phys.\ Rev.\ Lett.\ 94, 163202 (2005), doi:10.1103/PhysRevLett.94.163202.
\bibitem{jung06} R.~O.~Jung, T.~E.~Stone, J.~B.~Boffard, L.~W.~Anderson and C.~C.~Lin, Phys.\ Rev.\ A 73, 022722 (2006), doi:10.1103/PhysRevA.73.022722.
\bibitem{kannari85} F.~Kannari, M.~Obara and T.~Fujioka, J.\ Appl.\ Phys.\ 57, 4309 (1985), doi:10.1063/1.334590.
\bibitem{jacob76} J.~H.~Jacob and J.~A.~Mangano, Appl.\ Phys.\ Lett.\ 28, 724 (1976), doi:10.1063/1.88629.
\bibitem{vriens65} L.~Vriens, Physica 31, 385 (1965), doi:10.1016/0031-8914(65)90043-1.
\bibitem{gazeli21} K.~Gazeli et al., Phys.\ Plasmas 28, 043301 (2021), doi:10.1063/5.0041471.
\bibitem{gazeli21rev} K.~Gazeli, G.~Lombardi, X.~Aubert, C.~Y.~Duluard, S.~Prasanna and K.~Hassouni, Plasma 4, 145 (2021), doi:10.3390/plasma4010009.
\bibitem{chang78} R.~S.~F.~Chang and D.~W.~Setser, J.\ Chem.\ Phys.\ 69, 3885 (1978), doi:10.1063/1.437126.
\bibitem{dzierzega00} K.~Dzier\.z\k{e}ga, U.~Volz, G.~Nave and U.~Griesmann, Phys.\ Rev.\ A 62, 022505 (2000), doi:10.1103/PhysRevA.62.022505.
\bibitem{leichner74} P.~K.~Leichner and R.~J.~Ericson, Phys.\ Rev.\ A 9, 251 (1974), doi:10.1103/PhysRevA.9.251.
\bibitem{ueda84} K.~Ueda, S.~Kanada, A.~Sasaki and H.~Takuma, Rev.\ Laser Eng.\ 12, 357 (1984), doi:10.2184/lsj.12.357.
\bibitem{salamero79} Y.~Salamero et al., J.\ Phys.\ B 12, 419 (1979), doi:10.1088/0022-3700/12/3/019.
\bibitem{tracy76} C.~J.~Tracy and H.~J.~Oskam, J.\ Chem.\ Phys.\ 65, 1666 (1976), doi:10.1063/1.433312.
\bibitem{smith63} D.~S.~Smith and R.~Turner, Can.\ J.\ Phys.\ 41, 1949 (1963), doi:10.1139/p63-195.
\bibitem{quigley78} G.~P.~Quigley and W.~M.~Hughes, Appl.\ Phys.\ Lett.\ 32, 649 (1978), doi:10.1063/1.89882.
\bibitem{wieme81} W.~Wieme and J.~Lenaerts, J.\ Chem.\ Phys.\ 74, 483 (1981), doi:10.1063/1.440855.
\bibitem{morikawa89} E.~Morikawa et al., J.\ Chem.\ Phys.\ 91, 1469 (1989), doi:10.1063/1.457108.
\bibitem{bhattacharya70} A.~K.~Bhattacharya, Phys.\ Rev.\ A 1, 1196 (1970), doi:10.1103/PhysRevA.1.1196.
\bibitem{grossl81} M.~Gr\"ossl et al., J.\ Chem.\ Phys.\ 74, 1728 (1981), doi:10.1063/1.441261.
\bibitem{smith72} D.~Smith, A.~G.~Dean and I.~C.~Plumb, J.\ Phys.\ B 5, 2134 (1972), doi:10.1088/0022-3700/5/11/020.
\bibitem{perezrios18} J.~P\'erez-R\'ios and C.~H.~Greene, Phys.\ Rev.\ A 98, 062707 (2018), doi:10.1103/PhysRevA.98.062707.
\bibitem{shiu77} Y.-J.~Shiu and M.~A.~Biondi, Phys.\ Rev.\ A 16, 1817 (1977), doi:10.1103/PhysRevA.16.1817.
\bibitem{mikus78} O.~Mikus, J.\ Phys.\ D 11, L39 (1978), doi:10.1088/0022-3727/11/3/003.
\bibitem{mansbach69} P.~Mansbach and J.~Keck, Phys.\ Rev.\ 181, 275 (1969), doi:10.1103/PhysRev.181.275.
\bibitem{pack62} J.~L.~Pack, R.~E.~Voshall and A.~V.~Phelps, Phys.\ Rev.\ 127, 2084 (1962), doi:10.1103/PhysRev.127.2084.
\bibitem{dutton75} J.~Dutton, J.\ Phys.\ Chem.\ Ref.\ Data 4, 577 (1975), doi:10.1063/1.555525, Tables 1.9 and 7.4.
\bibitem{kruithof40} A.~A.~Kruithof, Physica 7, 519 (1940), doi:10.1016/S0031-8914(40)90043-X.
\bibitem{heylen71} A.~E.~D.~Heylen, Int.\ J.\ Electron.\ 31, 19 (1971), doi:10.1080/00207217108938196.
\bibitem{saito03} K.~Saito, S.~Sasaki, H.~Tawara, T.~Sanami and E.~Shibamura, IEEE Trans.\ Nucl.\ Sci.\ 50, 2452 (2003), doi:10.1109/TNS.2003.820615.
\bibitem{morozov08} A.~Morozov, T.~Heindl, R.~Kr\"ucken, A.~Ulrich and J.~Wieser, J.\ Appl.\ Phys.\ 103, 103301 (2008), doi:10.1063/1.2931000.
\bibitem{eckstrom88} D.~J.~Eckstrom et al., J.\ Appl.\ Phys.\ 64, 1691 (1988), doi:10.1063/1.341762.
\bibitem{hunter88} S.~R.~Hunter, J.~G.~Carter and L.~G.~Christophorou, Phys.\ Rev.\ A 38, 5539 (1988), doi:10.1103/PhysRevA.38.5539.
\bibitem{pack92} J.~L.~Pack, R.~E.~Voshall, A.~V.~Phelps and L.~E.~Kline, J.\ Appl.\ Phys.\ 71, 5363 (1992), doi:10.1063/1.350555.
\bibitem{nakamura90} Y.~Nakamura, in \emph{Nonequilibrium Effects in Ion and Electron Transport} (Springer, 1990), p.~363, doi:10.1007/978-1-4613-0661-0\_33.
\bibitem{dallarmi92} G.~Dall'Armi, K.~L.~Brown, P.~H.~Purdie and J.~Fletcher, Aust.\ J.\ Phys.\ 45, 185 (1992), doi:10.1071/PH920185.
\bibitem{bowe60} J.~C.~Bowe, Phys.\ Rev.\ 117, 1411 (1960), doi:10.1103/PhysRev.117.1411.
\bibitem{english53} W.~N.~English and G.~C.~Hanna, Can.\ J.\ Phys.\ 31, 768 (1953), doi:10.1139/p53-070.
\bibitem{kucukarpaci81} H.~N.~Kucukarpaci and J.~Lucas, J.\ Phys.\ D 14, 2001 (1981), doi:10.1088/0022-3727/14/11/008.
\bibitem{koizumi86} T.~Koizumi, E.~Shirakawa and I.~Ogawa, J.\ Phys.\ B 19, 2331 (1986), doi:10.1088/0022-3700/19/15/014.
\bibitem{alamin87} S.~A.~J.~Al-Amin and J.~Lucas, J.\ Phys.\ D 20, 1590 (1987), doi:10.1088/0022-3727/20/12/007.
\bibitem{bhattacharya79} A.~K.~Bhattacharya, J.\ Appl.\ Phys.\ 50, 6207 (1979), doi:10.1063/1.325754.
\bibitem{specht80} L.~T.~Specht, S.~A.~Lawton and T.~A.~DeTemple, J.\ Appl.\ Phys.\ 51, 166 (1980), doi:10.1063/1.327395.
\bibitem{goodman10} J.~Goodman and J.~Weare, Commun.\ Appl.\ Math.\ Comput.\ Sci.\ 5, 65 (2010), doi:10.2140/camcos.2010.5.65.
\bibitem{mckay79} M.~D.~McKay, R.~J.~Beckman and W.~J.~Conover, Technometrics 21, 239 (1979), doi:10.1080/00401706.1979.10489755.
\bibitem{stokes20} P.~W.~Stokes, D.~G.~Cocks, M.~J.~Brunger and R.~D.~White, Plasma Sources Sci.\ Technol.\ 29, 055009 (2020), doi:10.1088/1361-6595/ab85b6.
\bibitem{kingma14} D.~P.~Kingma and J.~Ba, arXiv:1412.6980 (2014), doi:10.48550/arXiv.1412.6980.
\bibitem{marchal18} F.~Marchal, N.~Sewraj, J.-P.~Gardou, N.~Merbahi and M.~Yousfi, ``Study of formation and decay of rare-gas excimers by laser-induced fluorescence'', in \emph{Photon Counting -- Fundamentals and Applications}, ed.\ N.~Britun and A.~Nikiforov (IntechOpen, 2018), doi:10.5772/intechopen.71942.
\bibitem{audouard88} E.~Audouard, P.~Laporte, J.-L.~Subtil and N.~Damany, J.\ Chem.\ Phys.\ 89, 6176 (1988), doi:10.1063/1.455434.
\bibitem{audouard91} E.~Audouard, P.~Laporte, J.~L.~Subtil and R.~Reininger, J.\ Chem.\ Phys.\ 95, 3283 (1991), doi:10.1063/1.460834.
\bibitem{turner67} R.~Turner, Phys.\ Rev.\ 158, 121 (1967), doi:10.1103/PhysRev.158.121.
\bibitem{larocque86} P.~E.~LaRocque, R.~H.~Lipson, P.~R.~Herman and B.~P.~Stoicheff, J.\ Chem.\ Phys.\ 84, 6627 (1986), doi:10.1063/1.450716.
\bibitem{killeen86} K.~P.~Killeen and J.~G.~Eden, J.\ Chem.\ Phys.\ 84, 6048 (1986), doi:10.1063/1.450794.
\bibitem{signorell01} R.~Signorell, U.~Hollenstein and F.~Merkt, J.\ Chem.\ Phys.\ 114, 9840 (2001), doi:10.1063/1.1370939.
\bibitem{cunningham72} A.~J.~Cunningham and R.~M.~Hobson, J.\ Phys.\ B 5, 1773 (1972), doi:10.1088/0022-3700/5/9/020.
\bibitem{firestone88} R.~F.~Firestone and W.~E.~Selander, J.\ Chem.\ Phys.\ 88, 717 (1988), doi:10.1063/1.454150.
\bibitem{oka74} T.~Oka, K.~V.~S.~Rama Rao, J.~L.~Redpath and R.~F.~Firestone, J.\ Chem.\ Phys.\ 61, 4740 (1974), doi:10.1063/1.1681797.
\bibitem{janssens87} H.~Janssens, M.~Vanmarcke, E.~Desoppere, R.~Bouciqu\'e and W.~Wieme, J.\ Chem.\ Phys.\ 86, 4935 (1987), doi:10.1063/1.452663.
\bibitem{mano22} R.~D.~P.~Mano, C.~A.~O.~Henriques, F.~D.~Amaro and C.~M.~B.~Monteiro, Phys.\ Lett.\ B 824, 136782 (2022), doi:10.1016/j.physletb.2021.136782.
\bibitem{kholin15} I.~V.~Kholin, A.~Yu.~L'dov, N.~N.~Ustinovskii and D.~A.~Zayarnyi, J.\ Russ.\ Laser Res.\ 36, 74 (2015), doi:10.1007/s10946-015-9479-y.
\bibitem{blagoev84} A.~B.~Blagoev, T.~M.~Mishonov and Tc.~K.~Popov, J.\ Phys.\ B 17, 435 (1984), doi:10.1088/0022-3700/17/3/013.
\bibitem{vriens80} L.~Vriens and A.~H.~M.~Smeets, Phys.\ Rev.\ A 22, 940 (1980), doi:10.1103/PhysRevA.22.940.
\bibitem{flannery81} M.~R.~Flannery, K.~J.~McCann and N.~W.~Winter, J.\ Phys.\ B 14, 3789 (1981), doi:10.1088/0022-3700/14/19/025.
\bibitem{blagoev81} A.~B.~Blagoev and N.~H.~Pilosov, Phys.\ Lett.\ A 83, 161 (1981), doi:10.1016/0375-9601(81)90875-6.
\bibitem{qiu18} J.~Qiu, Z.-W.~Cheng, X.-M.~Zhu and Y.-K.~Pu, J.\ Phys.\ D 51, 135202 (2018), doi:10.1088/1361-6463/aab0f7.
\bibitem{zeng05} J.~Zeng, J.~Wu, F.~Jin, G.~Zhao and J.~Yuan, Phys.\ Rev.\ A 72, 042707 (2005), doi:10.1103/PhysRevA.72.042707.
\bibitem{whitehead95} C.~A.~Whitehead, H.~Pournasr, M.~R.~Bruce, H.~Cai, J.~Kohel, W.~B.~Layne and J.~W.~Keto, J.\ Chem.\ Phys.\ 102, 1965 (1995), doi:10.1063/1.468763.
\bibitem{nistasd} A.~Kramida, Yu.~Ralchenko, J.~Reader and NIST ASD Team, NIST Atomic Spectra Database (version 5.12), National Institute of Standards and Technology (2024), doi:10.18434/T4W30F, read 15 October 2026.
\end{thebibliography}

\end{document}
